\documentclass[11pt]{article}

\usepackage[margin=1in]{geometry}
\usepackage{amsmath,amssymb}
\usepackage{booktabs}
\usepackage{longtable}
\usepackage{tabularx}
\usepackage{xcolor}
\usepackage{hyperref}

\hypersetup{
  colorlinks=true,
  linkcolor=blue,
  citecolor=blue,
  urlcolor=blue
}

\makeatletter
\renewcommand*\l@subsection{\@dottedtocline{2}{1.5em}{3.2em}}
\makeatother

\newcommand{\OPH}{\mathrm{OPH}}
\newcommand{\LCDM}{\Lambda\mathrm{CDM}}
\newcommand{\MOND}{\mathrm{MOND}}
\newcommand{\Msun}{M_{\odot}}
\newcommand{\aoph}{a_{0,\mathrm{OPH}}}
\newcommand{\aeff}{a_{0,\mathrm{eff}}}
\newcommand{\lambdac}{\lambda_{\mathrm{collar}}}
\newcommand{\Zsix}{\mathbb Z_6}
\newcommand{\dd}{\mathrm{d}}
\newcommand{\status}[1]{\textsf{#1}}

\title{Observer-Patch Holography and the Dark Matter Phenomenon}
\author{Bernhard Mueller}
\date{May 3, 2026}
\newcommand{\OPHPaperReleaseID}{r1520}
\newcommand{\OPHPaperReleaseDate}{July 9, 2026}
\newcommand{\OPHPaperReleaseBanner}{%
\begin{center}
\small\textbf{Paper release:} \texttt{\OPHPaperReleaseID}\quad\textbf{Released:} \OPHPaperReleaseDate
\end{center}
}


\begin{document}
\maketitle
\begin{center}
\small\textbf{Paper release:} \OPHPaperReleaseID\quad
\textbf{Released:} \OPHPaperReleaseDate
\end{center}

\begin{abstract}
In OPH, gravity is the stress bookkeeping of stable agreement between finite observer patches. Imperfect agreement leaves a record-repair remainder on overlap collars. That remainder is electromagnetically dark and gravitates as an information-defect stress correction.

The continuation gives a galaxy-limit equation, cluster and cosmology contracts, and a coefficient target. In old, settled, low-acceleration galaxies the correction can approximate the MOND acceleration law. In clusters and cosmology it must be treated as transported stress with its own equations. A scalar source table is not enough for CMB or growth predictions; the full likelihood evaluation requires a finite covariant packet parent and frozen source/solver/likelihood hashes.
\end{abstract}

\section*{What This Paper Contributes}

The dark matter evidence is not one observation. Galaxy rotation curves, the baryonic
Tully-Fisher relation, the radial acceleration relation, cluster lensing, CMB peaks, BAO, growth,
and \(S_8\) all ask for a gravitating sector that ordinary luminous matter does not supply. This
paper treats the MOND-like galaxy law as a clue about equilibrium response and the \(\Lambda\)CDM
evidence as a clue about transported stress.

The OPH move is to give that stress an observer-patch origin. Imperfect record repair leaves a
collar remainder. It has no luminous charge, so electromagnetic probes read no matter signal. It
gravitates because OPH puts stable collar remainders into the effective stress bookkeeping.
In settled low-acceleration galaxies the same mechanism yields the RAR-shaped response. In
clusters and cosmology it becomes a transported component with its own relaxation and likelihood
contracts. The paper is a continuation surface, so the joint fit remains the decisive gate.

\tableofcontents

\section{Plain Picture}

A galaxy does two jobs in this model.  Its baryons source the ordinary weak
field, and its finite observer-overlap collars carry a record-repair load.  The
load is dark to electromagnetic probes because ordinary luminous charges are
absent.  It gravitates because OPH puts carried collar remainders into the
effective stress bookkeeping.

The RAR shape has a simple origin.  The relevant repair opportunities live on
cuts, so their broad count scales like \(r_M/r=\sqrt{g_b/\aoph}\).  Independent
local opportunity counts give a Poisson zero-count law.  The active fraction is
\(1-\exp[-\lambdac\sqrt{g_b/\aoph}]\).  The conservative field equation
promotes that scalar response into a potential equation.  The simple algebraic
acceleration law is exact only when symmetry removes the curl term.

\section{Three Claim Layers}

The first layer is the static galaxy law.  It contains the RAR shape, the
BTFR limit, the effective dark density, the no-slip lensing rule, and the
\(\Zsix\)/Poisson coefficient.

The second layer is the dynamic dark component.  It treats the anomaly as a
transported stress that relaxes toward the static equilibrium over
\(\tau_{\mathrm{rec}}\).  Cluster offsets use this layer.

The third layer is cosmological dark matter.  It uses the same stress in
linear perturbation theory, with a background abundance, a relaxation rate,
and an equilibrium kernel.  CMB, BAO, lensing, growth, and \(S_8\) test this
layer.

\section{The Dark Matter Challenge}

In Newtonian gravity and in the weak-field limit of general relativity, a
test particle in a circular orbit satisfies
\begin{equation}
  \frac{v^2(r)}{r}=g(r)=\frac{G M(<r)}{r^2}.
\end{equation}
If almost all gravitating matter is visible baryonic matter, then outside the
bright disk the enclosed mass should stop growing quickly and the rotation
speed should fall as \(v(r)\propto r^{-1/2}\).  Disk galaxies do not behave
that way.  Their outer rotation curves often remain flat, so the inferred
dynamical mass grows approximately linearly with radius:
\begin{equation}
  M_{\mathrm{dyn}}(<r)=\frac{v^2 r}{G}.
\end{equation}

The galaxy problem is only one part of the evidence.  Clusters need more
gravity than the observed gas and galaxies provide.  Weak lensing maps in
merging systems such as the Bullet Cluster show gravitational potential peaks
separated from the X-ray gas \cite{clowe2006}.  The CMB and large-scale
structure prefer a gravitating component that is mostly non-baryonic, close to
pressureless, and present before galaxies settle \cite{planck2018}.  Planck's
base \(\LCDM\) parameter set gives
\(\Omega_c h^2=0.120\pm0.001\) and
\(\Omega_b h^2=0.0224\pm0.0001\), so most matter is not baryonic in that
model \cite{planck2018}.

The challenge includes more than ``flat rotation curves''.  A good theory has
to explain:
\begin{itemize}
\item galaxy rotation curves and the baryonic Tully-Fisher relation;
\item the tight radial acceleration relation, abbreviated RAR;
\item galaxy and cluster lensing;
\item cluster offsets in mergers;
\item CMB acoustic peaks, BAO, matter power, weak lensing, and \(S_8\).
\end{itemize}

\section{Standard Routes}

\subsection{Cold Dark Matter}

The standard route adds a new non-baryonic matter component.  In its simplest
form it is cold, effectively collisionless, and pressureless on cosmological
scales.  \(\LCDM\) uses the same component for CMB peaks, large-scale
structure, clusters, lensing, and galaxy halos.

The particle identity is absent.  Direct detection has produced no confirmed
dark matter particle.  Galaxy data show a tight
connection between baryonic acceleration and observed acceleration, much
tighter than a naive halo-plus-disk picture suggests.  Halo fitting also
introduces galaxy-level nuisance freedom, while the RAR looks like a law.

\subsection{MOND}

MOND should be treated as a clue, not as a failure.  Its basic phenomenology
starts from a critical acceleration \(a_0\), and in the deep low-acceleration
regime it gives
\begin{equation}
  g \simeq \sqrt{a_0 g_b}.
\end{equation}
For a point baryonic mass this implies
\begin{equation}
  v^4 = G M_b a_0,
\end{equation}
which is exactly the baryonic Tully-Fisher scaling \cite{milgrom1983}.  This
is MOND's great success: it goes straight at the galaxy regularity.

The pressure points start when the same static acceleration rule is asked to
do more jobs than a galaxy equilibrium law can naturally do:
\begin{itemize}
\item \emph{Relativistic completion and lensing.}  A nonrelativistic force law
does not by itself say which metric photons follow, or how the modified field
contributes to the stress side of Einstein's equations.  Relativistic
completions such as TeVeS exist \cite{bekenstein2004}, but the lensing sector
then depends on extra fields and closure assumptions.
\item \emph{Clusters.}  MOND reduces the cluster mass discrepancy but does not
remove the need for additional gravitating stress in clusters.  Merging systems
are especially sharp: if the anomaly is an instantaneous function of the local
baryonic acceleration, the lensing peak wants to remain too tightly tied to the
visible baryonic source.
\item \emph{The early universe.}  The CMB acoustic peaks, BAO, matter growth,
CMB lensing, and weak-lensing amplitudes need a homogeneous and perturbative
gravitating component before settled galaxies exist.  A static RAR law does
not by itself define that background abundance or its linear response.
\item \emph{Non-spherical and time-dependent systems.}  The algebraic
relation \(g\simeq\sqrt{a_0g_b}\) is clean in high-symmetry systems.  In real
three-dimensional fields, curl terms, environmental dependence, transport, and
relaxation can matter.
\end{itemize}

OPH does not treat the static RAR formula as the master law.  It approximates
MOND in the regime where the OPH stress sector has relaxed onto the old,
settled, low-acceleration galaxy branch.  It need not approximate MOND in
mergers, clusters, or the early universe, because those systems are not static
one-dimensional galaxy equilibria.  The underlying object is an effective
stress sector carried by finite overlap repair.  In old relaxed disks it
projects to a baryon-linked acceleration law.  In lensing it enters the
metric-potential bookkeeping.  In cluster mergers it is transported and
relaxes over a finite repair time.  In FLRW and perturbation theory it has to
be specified as a homogeneous charge, load current, and response kernel.

This is the sense in which the MOND pressure points dissolve in OPH.  OPH
treats those failures as wrong-regime extrapolations, without making the same
algebraic law more complicated.  The static galaxy formula is allowed to be a
very good equilibrium approximation, while lensing, clusters, and cosmology are
handled by the corresponding OPH stress, transport, and perturbation equations.
Those equations must pass their own falsifiers.  The claim is structural
separation, with no immunity from data.

\section{The OPH Route}

Observer-Patch Holography starts from finite observer patches and the demand
that neighboring patches agree on overlap data.  The relevant OPH inputs for
the dark branch are:

\begin{itemize}
\item finite overlap collars and fixed-cutoff collar algebras;
\item edge and cut registers on those collars;
\item read, compare, repair, commit, and refresh operations;
\item a finite carried modular remainder on the collar;
\item the recovered Einstein branch with a screen-capacity de Sitter scale.
\end{itemize}

The dark-sector proposal is:
\begin{quote}
The dark source is a transported modular/collar information-defect remainder.  It
gravitates because it contributes to the effective stress bookkeeping.  It is
electromagnetically dark because it is a record-repair remainder, not an
ordinary luminous matter species.
\end{quote}
No new particle species is postulated.  The effective dark stress is the
gravitating part of finite-observer information repair.
\begin{equation}
  G_{ab}+\Lambda g_{ab}
  =
  8\pi G
  \left(
    T_{ab}^{\mathrm{SM}}
    +T_{ab}^{\mathcal I}
    +T_{ab}^{\mathrm{repair}}
  \right).
  \label{eq:informationstresssurface}
\end{equation}
The symbol \(A\) in formulas below denotes this anomaly or information-defect
component.  It is an effective stress sector, with no literal particle species
claim.

This differs from particle dark matter and MOND in a useful way.  It is
stress-first like dark matter, so lensing and cosmology have a natural slot.
It is baryon-linked like MOND only on the settled galaxy branch, because that
equilibrium response is activated by collar/cut repair opportunities sourced by
baryonic acceleration.  The important contrast with MOND is not the
square-root galaxy law; OPH intentionally reproduces that approximation where
the equilibrium assumptions apply.  The contrast is where the law lives.  In
MOND, the galaxy formula is often treated as the modification itself.  In OPH
it is the equilibrium projection of a transported stress component.  That is
why lensing, clusters, and cosmology can have their own equations without
giving up the galaxy regularity.

\section{Static Galaxy Derivation}

\subsection{Source From The Carried Collar Remainder}

On the support-visible BW branch the primary theorem is the cap automorphism identity. In the
special type-I or effective local representation where the geometric generator may be written
inside the cap algebra, the OPH modular decomposition has the local form
\begin{equation}
  K_C = 2\pi B_C + K_C^{(\mathrm{anom})}+\mathrm{const}.
\end{equation}
At finite regulator stage, the nonadditive modular part is carried on the
overlap collar.  The dark branch packages the galaxy-scale positive remainder
as
\begin{equation}
  R_C := I(A:D|B)\ge 0,
\end{equation}
where \(I(A:D|B)\) is conditional mutual information.  This is a
phenomenological carried-defect variable on the continuation branch: small CMI
supplies a recovered comparison state and a fixed-collar replacement modulus,
not a one-shot exact Markov normal form or an exact finite-cutoff Einstein
source. The corresponding
effective rest-energy density is written
\begin{equation}
  \rho_A c^2
    = \frac{15}{8\pi^2\ell^4} R_C.
  \label{eq:rhoA}
\end{equation}
The weak-field Poisson equation becomes
\begin{equation}
  \nabla^2 \Phi = 4\pi G(\rho_b+\rho_A).
  \label{eq:poisson}
\end{equation}

\paragraph{Conditional theorem 1, attractive sign.}
If the carried galaxy remainder is \(R_C=I(A:D|B)\) and the coarse-grained
domain measure is nonnegative, then the anomaly is attractive in the
Newtonian weak-field equation.

\paragraph{Proof.}
Conditional mutual information is nonnegative, so \(R_C\ge0\).  Equation
\eqref{eq:rhoA} gives \(\rho_A\ge0\).  Equation \eqref{eq:poisson} then shows
that the anomaly enters with the same attractive sign as baryonic mass.
\(\square\)

The source identification is a theorem target.

\subsection{The Acceleration Variable}

Define the baryonic acceleration ratio
\begin{equation}
  x := \frac{g_b}{\aoph}.
\end{equation}
For a point baryonic source,
\begin{equation}
  g_b(r)=\frac{G M_b}{r^2},
  \qquad
  r_M := \sqrt{\frac{G M_b}{\aoph}},
\end{equation}
so
\begin{equation}
  x=\left(\frac{r_M}{r}\right)^2,
  \qquad
  \sqrt{x}=\frac{r_M}{r}.
  \label{eq:sqrtx}
\end{equation}

The OPH support assumption is that settled-galaxy repair opportunities live
on codimension-one collar/cut support.  A screen-area count would scale as
\(x\).  A cut-count scales as \(\sqrt{x}\), as in \eqref{eq:sqrtx}.

\subsection{Poisson Activation}

Let the mean number of independent local repair opportunities be
\begin{equation}
  \mu(x)=\lambdac\sqrt{x}.
\end{equation}
The Poisson law does not follow from fixing this mean alone.
Fixed mean alone would select a geometric distribution on nonnegative
integers.  The OPH claim needs the stronger collar-count premise: independent
increments on disjoint collar intervals, rare local repair events, and
refinement stability.  Partition an interval \(I\) with mean \(\mu(I)\) into
\(m\) equal subintervals.  If each subinterval has one activation opportunity
with probability \(\mu(I)/m+O(m^{-2})\), and the probability of two or more
events in one subinterval is \(O(m^{-2})\), then for every fixed \(n\),
\begin{equation}
  \Pr[N(I)=n]
  =
  \lim_{m\to\infty}
  {m\choose n}
  \left(\frac{\mu(I)}{m}\right)^n
  \left(1-\frac{\mu(I)}{m}\right)^{m-n}
  =
  e^{-\mu(I)}\frac{\mu(I)^n}{n!}.
  \label{eq:poissonlaw}
\end{equation}
In particular, the zero-count probability is
\begin{equation}
  \Pr[N(I)=0]
  =
  \lim_{m\to\infty}
  \left(1-\frac{\mu(I)}{m}\right)^m
  =
  e^{-\mu(I)}.
\end{equation}
The probability that at least one repair opportunity is active is
\begin{equation}
  p(x)=1-\Pr[N=0]
      =1-\exp[-\lambdac\sqrt{x}].
  \label{eq:activation}
\end{equation}

The activation count is only the counting part of the static law.  The source
closure says that active scalar repair opportunities recover baryonic flux in
the settled one-dimensional branch:
\begin{equation}
  g_b = p(x) g_{\mathrm{obs}}.
  \label{eq:fluxrecovery}
\end{equation}
Thus
\begin{equation}
  g_{\mathrm{obs}}
    = \nu_{\OPH}(x) g_b,
  \qquad
  \nu_{\OPH}(x)
    = \frac{1}{1-\exp[-\lambdac\sqrt{x}]}.
  \label{eq:ophrar}
\end{equation}

\paragraph{Conditional theorem 2, activation law.}
Assume codimension-one collar/cut support, independent increments on disjoint
collar intervals, rare local repair events, refinement stability, and the
settled scalar flux-recovery closure \eqref{eq:fluxrecovery}.  Then the
settled one-dimensional galaxy response is \eqref{eq:ophrar}.

\paragraph{Proof.}
Codimension-one support gives a mean count proportional to \(\sqrt{x}\).
The scalar coefficient is \(\lambdac\).  Independent increments and rare-event
refinement give the Poisson law \eqref{eq:poissonlaw}.  The active probability is
the complement of the zero-count event, which is \eqref{eq:activation}.
Solving \(g_b=p(x)g_{\mathrm{obs}}\) in one-dimensional symmetry gives
\eqref{eq:ophrar}. \(\square\)

\subsection{Conservative Potential Equation}

For general disks and three-dimensional baryonic profiles, the algebraic vector
rule
\begin{equation}
  \mathbf g_{\mathrm{alg}}
  =
  \nu_{\OPH}(|\mathbf g_b|/\aoph)\mathbf g_b
\end{equation}
need not be curl-free because
\begin{equation}
  \nabla\times\mathbf g_{\mathrm{alg}}
  =
  \nabla\nu_{\OPH}\times\mathbf g_b.
\end{equation}
A gravitational acceleration must be derived from a potential.  The primary
static OPH equation is therefore
\begin{equation}
  \nabla^2\Phi
  =
  \nabla\cdot
  \left[
    \nu_{\OPH}\!\left(\frac{|\nabla\Phi_b|}{\aoph}\right)
    \nabla\Phi_b
  \right],
  \qquad
  \mathbf g_{\mathrm{obs}}=-\nabla\Phi,
  \label{eq:ophpotential}
\end{equation}
with
\begin{equation}
  \nabla^2\Phi_b=4\pi G\rho_b.
\end{equation}
In one-dimensional symmetry, \eqref{eq:ophpotential} integrates to
\begin{equation}
  g_{\mathrm{obs}}
  =
  \nu_{\OPH}(g_b/\aoph)g_b+C.
\end{equation}
The boundary condition \(g_{\mathrm{obs}}\to0\) when the enclosed baryonic
flux vanishes sets \(C=0\), so \eqref{eq:ophrar} is exact for spherical,
planar, or one-dimensional symmetry.  The SPARC tables in this note use the
algebraic diagnostic; a full disk-potential solver is required for
publication-grade rotation and lensing claims.

\subsection{Newtonian And Deep-IR Limits}

For \(x\gg1\), \(p(x)\to1\), hence
\begin{equation}
  g_{\mathrm{obs}}\to g_b.
\end{equation}
For \(x\ll1\),
\begin{equation}
  p(x)=\lambdac\sqrt{x}+O(x),
\end{equation}
so
\begin{equation}
  g_{\mathrm{obs}}
  \simeq \frac{g_b}{\lambdac\sqrt{g_b/\aoph}}
  = \sqrt{\frac{\aoph}{\lambdac^2}g_b}
  = \sqrt{\aeff g_b},
\end{equation}
where
\begin{equation}
  \aeff = \frac{\aoph}{\lambdac^2}.
  \label{eq:aeff}
\end{equation}
For a point source,
\begin{equation}
  v^4 = G M_b \aeff.
  \label{eq:btfr}
\end{equation}

\paragraph{Theorem 3, BTFR.}
The OPH activation law implies the baryonic Tully-Fisher scaling in the
deep-IR limit.

\paragraph{Proof.}
For a circular orbit, \(v^2/r=g_{\mathrm{obs}}\).  In the deep-IR limit,
\(g_{\mathrm{obs}}=\sqrt{\aeff G M_b/r^2}\).  Multiplying by \(r\) gives
\(v^2=\sqrt{G M_b\aeff}\), hence \eqref{eq:btfr}. \(\square\)

\subsection{Effective Dark Density}

The anomaly can be represented as an equilibrium effective density:
\begin{equation}
  \mathbf g_A=(\nu_{\OPH}-1)\mathbf g_b,
  \qquad
  \rho_{A,\mathrm{eq}}
   =-\frac{1}{4\pi G}\nabla\cdot[(\nu_{\OPH}-1)\mathbf g_b].
  \label{eq:rhoaeq}
\end{equation}
For a spherical system,
\begin{equation}
  M_{\mathrm{dyn}}(r)=\nu_{\OPH}(x) M_b(r),
  \qquad
  M_A(r)=[\nu_{\OPH}(x)-1]M_b(r).
\end{equation}
In the deep-IR point-source limit, \(M_{\mathrm{dyn}}(r)\propto r\), which is
the effective mass profile behind flat rotation curves.

For a point baryonic source, the exact equilibrium expression is
\begin{equation}
  M_A(r)
  =
  \frac{M_b}{\exp(\lambdac r_M/r)-1},
\end{equation}
and
\begin{equation}
  \rho_A(r)
  =
  \frac{
    M_b\lambdac r_M\exp(\lambdac r_M/r)
  }{
    4\pi r^4[\exp(\lambdac r_M/r)-1]^2
  }.
\end{equation}
For an extended spherical baryonic profile, positivity of the settled
effective density is the monotonicity condition
\begin{equation}
  \frac{\dd}{\dd r}
  \left(
    [\nu_{\OPH}(x(r))-1]M_b(r)
  \right)
  \ge 0.
\end{equation}
Profiles that violate this condition need finite pressure, shear, or transport
terms instead of the settled dust equilibrium alone.

\section{The Finite-Collar Coefficient}

\subsection{Unit Branch}

The default branch is
\begin{equation}
  \lambdac=1,
  \qquad
  \aeff=\aoph.
\end{equation}
This branch sets the finite-collar opportunity density to one.

\subsection{\(\Zsix\) Reserve Branch}

The OPH gauge structure contains the realized quotient
\begin{equation}
  G_{\mathrm{phys}}
    = \frac{SU(3)\times SU(2)\times U(1)}{\Zsix}.
\end{equation}
The protected shared-edge reserve has mean
\begin{equation}
  \epsilon_{\Zsix}
    = \frac{\bar\ell_{\mathrm{shared}}}{|\Zsix|}
    = \frac{P/4}{6}
    = \frac{P}{24}.
\end{equation}
Using
\begin{equation}
  P=1.630968209403959,
\end{equation}
gives
\begin{equation}
  \epsilon_{\Zsix}=0.06795700872516496.
\end{equation}

\paragraph{Conditional bridge 4, quotient-edge co-registration.}
Galaxy scalar repair opportunities are quotient-edge cut-register events on
the same edge-center resolution that carries the protected \(\Zsix\) reserve.
The \(\Zsix\) reserve therefore co-registers with the scalar activation count.

\paragraph{Conditional argument.}
The scalar source is a quotient-local recoverability defect, so pure
\(\Zsix\) center motion changes no realized matter or Higgs state and is
inactive for scalar sourcing.  If the scalar count and the protected reserve
are forced onto the same center-projector resolution, the reserve removes
active scalar opportunities from that same layer. \(\square\)

\paragraph{Proof obligation.}
The fixed-cutoff screen-net construction must prove that galaxy scalar
activation cannot live in a separate bulk, screen-area, or noncentral collar
channel.

\paragraph{Conditional bridge 5, local Poisson reserve survival.}
If the reserve occupancy in a co-registered local scalar slot at transverse
collar coordinate \(y\) is Poisson with mean \(\epsilon_{\Zsix}(y)\), then the
local scalar survival factor is
\begin{equation}
  \lambda_{\mathrm{slot}}(y)
  =
  \Pr[N_{\Zsix}(y)=0]
  =
  \exp[-\epsilon_{\Zsix}(y)].
  \label{eq:z6localsurvival}
\end{equation}

\paragraph{Proof.}
A Poisson variable with mean \(\epsilon_{\Zsix}\) has zero-count probability
\(\exp(-\epsilon_{\Zsix})\).  A scalar opportunity survives only when its
co-registered reserve count is zero.  This proves the local survival factor.
\(\square\)

\subsection{Finite Thickness}

Let \(y\) be the transverse collar coordinate and \(w(y)\) the normalized
activation weight:
\begin{equation}
  \int \dd y\,w(y)=1.
\end{equation}
The finite-thickness branch gives
\begin{equation}
  \lambdac
   =\int \dd y\, w(y)\exp[-\epsilon_{\Zsix}(y)].
  \label{eq:thick}
\end{equation}
If
\begin{equation}
  \int \dd y\,w(y)\epsilon_{\Zsix}(y)=P/24,
\end{equation}
then Jensen's inequality gives
\begin{equation}
  \lambdac \ge \exp(-P/24).
\end{equation}
Equality holds when \(\epsilon_{\Zsix}(y)\) is constant on the weighted
collar.

\paragraph{Conditional bridge 6, exact finite-thickness coefficient.}
On the product transverse regulator with one quotient trace, the exact
finite-thickness coefficient is \(\lambdac=\exp(-P/24)\).

\paragraph{Conditional argument.}
If the finite-thickness reserve density is uniform on the scalar-weighted
collar layer, then \(\epsilon_{\Zsix}(y)=\epsilon_{\Zsix}\).  Equation
\eqref{eq:thick}
becomes
\begin{equation}
  \lambdac
  =\exp(-\epsilon_{\Zsix})\int \dd y\,w(y)
  =\exp(-\epsilon_{\Zsix}).
\end{equation}
Using \(\epsilon_{\Zsix}=P/24\) proves the claim. \(\square\)

On that exact branch,
\begin{equation}
  \lambdac=\exp(-P/24)=0.934300639489386\ldots
  \label{eq:z6lambda}
\end{equation}
and by \eqref{eq:aeff}
\begin{equation}
  \aeff=1.179018696\times10^{-10}\,\mathrm{m\,s^{-2}}.
\end{equation}

\paragraph{Failure consequence.}
Without finite-thickness uniformity, the coefficient is the integral
\eqref{eq:thick}.  Exact \(\exp(-P/24)\) is not available as a theorem-grade
number.

\section{Lensing And Relativistic Stress}

A density source by itself fixes the Newtonian Poisson equation.  Lensing also
depends on the spatial stress.  The OPH no-slip claim uses the transported
record-number stress and its spatial moments.

The conditional OPH stress closure is a kinetic record measure.  At each
coarse-grained event, the finite collar packet emits a positive measure
\(\dd\mathcal E_x(v)\) on the future unit timelike shell.  Define
\begin{equation}
  J_{\mathcal I}^a(x)=\int v^a\,\dd\mathcal E_x(v),
  \qquad
  T_{\mathcal I}^{ab}(x)=\int v^a v^b\,\dd\mathcal E_x(v).
  \label{eq:kineticstress}
\end{equation}
This stress is symmetric and future-positive.  Relative to an observer
\(u^a\), it decomposes as
\begin{equation}
  T_{\mathcal I}^{ab}
  =
  \rho_{\mathcal I}u^a u^b
  +P_{\mathcal I}h^{ab}
  +q_{\mathcal I}^{(a}u^{b)}
  +\pi_{\mathcal I}^{ab}.
  \label{eq:stressdecomposition}
\end{equation}
The settled galaxy branch is the monokinetic rest branch
\(\dd\mathcal E_x(v)=\rho_{\mathcal I}(x)\delta(v-u_{\mathcal I})\dd\Omega\).
Then
\begin{equation}
  T_{\mathcal I}^{ab}
  =
  \rho_{\mathcal I}u_{\mathcal I}^a u_{\mathcal I}^b,
  \qquad
  P_{\mathcal I}=0,
  \qquad
  q_{\mathcal I}^a=0,
  \qquad
  \pi_{\mathcal I}^{ab}=0.
  \label{eq:duststress}
\end{equation}
In Newtonian gauge,
\begin{equation}
  \dd s^2
  =a^2[-(1+2\Psi)\dd\eta^2+(1-2\Phi)\dd x^2].
\end{equation}
The off-diagonal spatial Einstein equation gives
\begin{equation}
  k^2(\Phi-\Psi)
  =12\pi G a^2\sum_i(\rho_i+P_i)\sigma_i.
\end{equation}
The \(\mathcal I\)-sector contribution to the right-hand side vanishes on
\eqref{eq:duststress}.  If baryonic, neutrino, radiation, and repair-sector
anisotropic stresses are absent or handled separately, then
\begin{equation}
  \Phi=\Psi.
\end{equation}
The lensing potential is therefore the same potential that controls dynamics:
\begin{equation}
  \Phi_{\mathrm{lens}}=\frac{\Phi+\Psi}{2}=\Phi.
\end{equation}

This is where OPH behaves differently from a force-only MOND law.  The
correction is attached to stress bookkeeping, so the same potential controls
slow matter and lensing on settled no-slip galaxy states.  Cluster mergers and
cosmological perturbations use the full decomposition
\eqref{eq:stressdecomposition}; a velocity-spread branch can carry pressure,
heat flux, or anisotropic stress.

\section{Clusters And Transport}

The static galaxy law is an equilibrium law.  Merging clusters need a
transported information-defect state.  In the cold transported branch, the
kinetic stress \eqref{eq:kineticstress} reduces to a density and velocity
field:
\begin{align}
  \partial_t\rho_{\mathcal I}
    +\nabla\cdot(\rho_{\mathcal I}\mathbf v_{\mathcal I})
    &=-\frac{\rho_{\mathcal I}
      -\rho_{\mathcal I,\mathrm{eq}}[\rho_b]}{\tau_{\mathrm{rec}}},
    \label{eq:transport1}\\
  \partial_t\mathbf v_{\mathcal I}
    +(\mathbf v_{\mathcal I}\cdot\nabla)\mathbf v_{\mathcal I}
    &=-\nabla\Phi,
    \label{eq:transport2}\\
  \nabla^2\Phi&=4\pi G(\rho_b+\rho_{\mathcal I}).
    \label{eq:transport3}
\end{align}
Velocity-spread branches use the moment equations of
\eqref{eq:stressdecomposition} instead of \eqref{eq:transport1}.

The repair-commit theorem gives the form of the relaxation term.  Once the
finite collar packet algebra is declared, the fixed-cutoff packet state space,
proposal rule, equilibrium weights, reversible transition matrix, and spectral
gap are determined by packet microphysics.  Let \(\gamma_{\mathrm{rec}}\) be
that dimensionless gap and \(t_{\mathrm{commit}}\) the physical commit time.
Then
\begin{equation}
  \Gamma_{\mathrm{rec}}
   =\frac{\gamma_{\mathrm{rec}}}{t_{\mathrm{commit}}},
  \qquad
  \tau_{\mathrm{rec}}
   =\frac{t_{\mathrm{commit}}}{\gamma_{\mathrm{rec}}}.
\end{equation}

Bianchi consistency requires a repair-exchange stress:
\begin{equation}
  \nabla_a(T_b^{ab}+T_{\mathcal I}^{ab}+T_R^{ab}+\cdots)=0.
\end{equation}
In quasi-covariant notation,
\begin{equation}
  \nabla_a T_{\mathcal I}^{ab}
  =
  -\Gamma_{\mathrm{rec}}
  (\rho_{\mathcal I}-\rho_{\mathcal I,\mathrm{eq}})
  u_{\mathcal I}^b,
\end{equation}
and \(T_R^{ab}\) carries the opposite exchange flux.

An OPH-side no-fit clock candidate is the Jacobi mismatch rate.  Let
\(E_{ij}=C_{0i0j}\) be the electric Weyl tidal tensor measured in the local
free-fall frame, equivalently the trace-free Newtonian Hessian in the weak
field.  The homogeneous FLRW Ricci focusing term is subtracted, because it
does not create a local overlap-visible cluster offset.  The scalar rate
\begin{equation}
  \Gamma_J
  =
  \left(\frac{E_{ij}E^{ij}}{6}\right)^{1/4},
  \qquad
  \tau_J=\Gamma_J^{-1}
  \label{eq:jacobiclock}
\end{equation}
is invariant under local rotations and has the right units.  For a Newtonian
point source,
\begin{equation}
  E_{ij}E^{ij}=6\left(\frac{GM}{r^3}\right)^2,
  \qquad
  \tau_J=\sqrt{\frac{r^3}{GM}}.
\end{equation}
The numerical scale is:
\begin{center}
\begin{tabular}{@{}lrrr@{}}
\toprule
Scale & \(M/\Msun\) & \(r\) in kpc & \(\tau_J\) in Gyr \\
\midrule
Inner spiral disk & \(6.0\times10^{10}\) & \(10\) & \(0.0608673\) \\
Outer spiral disk & \(6.0\times10^{10}\) & \(30\) & \(0.316276\) \\
Dwarf scale & \(1.0\times10^9\) & \(5\) & \(0.166692\) \\
Cluster core & \(1.0\times10^{14}\) & \(300\) & \(0.244986\) \\
Massive cluster aperture & \(1.0\times10^{15}\) & \(1000\) & \(0.471476\) \\
\bottomrule
\end{tabular}
\end{center}
For \(d_0=200\,\mathrm{kpc}\) and \(t=0.2\,\mathrm{Gyr}\), the massive-cluster
aperture row gives \(d(t)=130.859\,\mathrm{kpc}\).

The conditional physical-time conversion is
\begin{equation}
  \Gamma_{\mathrm{rec}}(x,t)=\Gamma_J(x,t),
  \qquad
  t_{\mathrm{commit}}(x,t)
  =\frac{\gamma_{\mathrm{rec}}(x,t)}{\Gamma_J(x,t)} .
\end{equation}
The static RAR law does not determine this clock.  Any positive relaxation
rate leaves the same equilibrium density.  The Jacobi clock requires the
premise that repair timing is set by the overlap-visible tidal mismatch rate
of transported information-defect records.  The theorem-grade packet clock
must derive this premise without using cluster offsets as selectors.

For a one-mode merger offset,
\begin{equation}
  d(t)=d_0\exp(-t/\tau_{\mathrm{rec}}),
  \qquad
  \tau_{\mathrm{rec}}=\frac{t}{\ln(d_0/d(t))}.
\end{equation}
With \(d_0=200\,\mathrm{kpc}\) and \(t=0.2\,\mathrm{Gyr}\):
\begin{center}
\begin{tabular}{@{}rr@{}}
\toprule
Retained offset & Required \(\tau_{\mathrm{rec}}\) \\
\midrule
\(50\,\mathrm{kpc}\) & \(0.144\,\mathrm{Gyr}\) \\
\(100\,\mathrm{kpc}\) & \(0.288\,\mathrm{Gyr}\) \\
\(134\,\mathrm{kpc}\) & \(0.500\,\mathrm{Gyr}\) \\
\(150\,\mathrm{kpc}\) & \(0.696\,\mathrm{Gyr}\) \\
\(164\,\mathrm{kpc}\) & \(1.000\,\mathrm{Gyr}\) \\
\bottomrule
\end{tabular}
\end{center}
Settled galaxies impose
\begin{equation}
  \tau_{\mathrm{rec}}<\frac{T_{\mathrm{set}}}{\ln(1/\epsilon)}
\end{equation}
if the allowed tracking error after settling time \(T_{\mathrm{set}}\) is
\(\epsilon\).  For \(T_{\mathrm{set}}=5\,\mathrm{Gyr}\), the bounds are
\(2.17\,\mathrm{Gyr}\), \(1.67\,\mathrm{Gyr}\), and \(1.09\,\mathrm{Gyr}\)
for \(10\%\), \(5\%\), and \(1\%\) tracking error.

\section{CMB, BAO, Growth, And \(S_8\)}

The linear cosmology contract has two distinct branches.  A conserved
homogeneous information-defect branch obeys
\begin{equation}
  \rho_A' + 3\mathcal H\rho_A=0,
  \qquad
  \rho_A(a)=\rho_{A0}a^{-3},
  \label{eq:conservedbranch}
\end{equation}
where primes denote conformal time and \(\mathcal H=a'/a\).  This branch
behaves like a cold stress component if its kinetic measure is monokinetic and
its homogeneous charge \(Q_A=a^3\rho_A V_{\mathrm{com}}\) is selected by OPH
screen data.

\paragraph{No-go boundary for scalar parents.}
A scalar parent that emits only
\(\rho_A(a)\), \(\rho_{A,\mathrm{eq}}(a)\), and \(B_A(k,a)\) is a source
candidate, not a physical Boltzmann closure.  It does not determine
\[
  w_A(a),\quad c_{s,A}^2(k,a),\quad \sigma_A(k,a),\quad Q_A^\mu,\quad
  \Gamma_{\mathrm{rec}}(k,a),
\]
and it does not by itself fix gauge-independent perturbations or the recipient
stress needed when repair transfers energy-momentum.  In the physical source
contract \(B_A\) is built from the anomaly-frame baryon density
\[
  n_b^{(A)}=-u_{A\mu}J_b^\mu ,
\]
not from a gauge-fixed coordinate density.  The physical object is a
finite covariant collar-packet parent
\[
  \mathcal P_r[X,g]=(C_r,Z_r,A_r,R_r,G_r,\pi_r[X],L_r[X],Q_r,D_r),
\]
where \(Z_r\) is a finite packet-state set split into anomaly and recipient
channels, \(\pi_r\) is the convex equilibrium packet functional, \(L_r\) is the
repair generator, \(Q_r\) records reaction channels, and \(D_r\) records any
causal auxiliary response.  The parent must emit the stress moments and source
projections consumed by the Boltzmann solver; otherwise CMB/BAO/growth rows
remain diagnostics.

For the source-only primordial bridge, this closed branch has the stricter
contract
\begin{equation}
  \nabla_aJ_A^a=0,\qquad
  T_A^{ab}=\rho_Au^au^b,\qquad
  Q_A^a=0 .
  \label{eq:closed-anomaly-primordial-contract}
\end{equation}
It may enter the rank-one primordial theorem only when \(\rho_A\) is a
function of the same source clock as the rest of the total stress.  Transport
of a supplied \(Q_A\) is not an OPH derivation of the homogeneous abundance.

A repair-exchange branch instead relaxes toward a finite-collar equilibrium
source:
\begin{equation}
  \rho_A' + 3\mathcal H\rho_A
  = -a\Gamma_{\mathrm{rec}}(\rho_A-\rho_{A,\mathrm{eq}}).
  \label{eq:repairbackground}
\end{equation}
The factor \(a\Gamma_{\mathrm{rec}}\) is the conformal-time rate associated
with the physical repair rate.  Define
\begin{equation}
  q_A=\frac{\rho_{A,\mathrm{eq}}}{\rho_A},
  \qquad
  \delta_{A,\mathrm{eq}}(k,a)=B_A(k,a)\delta_b(k,a).
\end{equation}
For a pressureless no-slip repair branch with energy exchange parallel to
the anomaly four-velocity, the Newtonian-gauge perturbations reduce to
\begin{align}
  \delta_A'
    &=-\theta_A+3\Phi'
      -a\Gamma_{\mathrm{rec}}q_A
       (\delta_A-\delta_{A,\mathrm{eq}}),
       \label{eq:deltaA}\\
  \theta_A'
    &=-\mathcal H\theta_A+k^2\Psi.
       \label{eq:thetaA}
\end{align}
The no-slip statement here concerns the information-defect sector.  The total
metric slip depends on baryons, radiation, neutrinos, repair stress, and any
anisotropic stress in a velocity-spread branch.

For the same primordial bridge, the repair-exchange branch is admissible only
with an explicit transfer receipt:
\begin{equation}
  \nabla_aT_A^{ab}=Q_A^b,\qquad
  \Xi_a^A=0,\qquad
  h^a{}_bQ_A^b=0
  \label{eq:repair-anomaly-primordial-contract}
\end{equation}
or with interval bounds on the two defects.  Current dark-sector work does not
yet derive the homogeneous anomaly amplitude \(Q_A\), and the finite-collar
kinetic stress construction remains conditional.  Therefore the first
source-only primordial certification run should use
\[
  \texttt{dark\_continuation = OFF}.
\]
A supplied dark abundance may be used only as
\[
  \texttt{dark\_continuation = CONDITIONAL\_SOURCE\_STATE},
\]
not as a derived OPH primordial prediction.

If \(\Gamma_{\mathrm{rec}}>0\), the recipient sector is not optional.  The
finite parent must expose a recipient stress \(T_R^{ab}\) and reaction channel
such that
\[
  \nabla_a(T_A^{ab}+T_R^{ab}+T_{\mathrm{other}}^{ab})=0 .
\]
Without that explicit stress, a repair sink/source can mimic a scalar damping
term but cannot be promoted to an Einstein-source prediction.

The finite scalar-load repair diagnostic gives
\[
  \gamma_{\mathrm{rec}}=0.050717471,
  \qquad
  \tau_{\mathrm{rec}}/t_{\mathrm{commit}}=19.717071 .
\]
It maps a cluster-scale relaxation window to the following commit times:
\begin{center}
\begin{tabular}{@{}rr@{}}
\toprule
\(\tau_{\mathrm{rec}}\) & \(t_{\mathrm{commit}}\) \\
\midrule
\(0.3\,\mathrm{Gyr}\) & \(15.2152\,\mathrm{Myr}\) \\
\(0.5\,\mathrm{Gyr}\) & \(25.3587\,\mathrm{Myr}\) \\
\(1.0\,\mathrm{Gyr}\) & \(50.7175\,\mathrm{Myr}\) \\
\bottomrule
\end{tabular}
\end{center}
For the same relaxation window, a simple background timing diagnostic gives:
\begin{center}
\begin{tabular}{@{}rrrr@{}}
\toprule
\(\tau_{\mathrm{rec}}\) & \(z=0\) & \(z=10\) & \(z=1100\) \\
\midrule
\(0.3\,\mathrm{Gyr}\) & \(48.4092\) & \(2.36213\) & \(0.00205644\) \\
\(0.5\,\mathrm{Gyr}\) & \(29.0455\) & \(1.41728\) & \(0.00123387\) \\
\(1.0\,\mathrm{Gyr}\) & \(14.5228\) & \(0.708639\) & \(0.000616933\) \\
\bottomrule
\end{tabular}
\end{center}
For a hybrid branch, repair relaxation is tiny during acoustic physics and
large in low-Hubble settled environments.  A cosmology claim requires a
finite covariant source artifact, Boltzmann implementation, frozen likelihood
hashes, and likelihood tests against CMB TT/TE/EE, CMB lensing, BAO, weak
lensing, redshift-space distortions, and \(S_8\).

\section{Abundance And Kernel Targets}

The static galaxy law fixes a settled response around resolved baryonic
inhomogeneity.  A cosmological dark-matter branch needs OPH source outputs:
\begin{equation}
  \rho_A(a),
  \qquad
  B_A(k,a),
  \qquad
  \Gamma_{\mathrm{rec}}(k,a),
  \qquad
  w_A,\ c_{s,A}^2,\ \sigma_A,\ Q_A^\mu .
\end{equation}
The first is the homogeneous anomaly abundance.  The second is the linear
equilibrium-source kernel in
\begin{equation}
  \delta_{A,\mathrm{eq}}(k,a)=B_A(k,a)\delta_b(k,a).
\end{equation}
The remaining entries close the stress tensor and exchange law.

\subsection{Static FLRW Boundary}

The nonlinear galaxy law cannot be inserted directly into FLRW perturbation
theory.  In an exactly homogeneous background there is no preferred baryonic
acceleration vector.  If one tries to Taylor-expand the static RAR law at
\(\mathbf g_b=0\), the susceptibility is singular because
\(\nu_{\OPH}(x)\sim(\lambdac\sqrt{x})^{-1}\).  The static RAR branch therefore
does not select a homogeneous anomaly density or a universal CMB kernel.

This is a useful boundary.  It prevents a hidden fit from being smuggled into
the galaxy law.  The CMB branch has to declare its background anomaly charge
and its perturbation kernel from OPH state selection or finite-collar
microphysics.

\subsection{Flat Capacity-Saturated State Selection}

One clean homogeneous branch is flat capacity saturation.  The de Sitter screen
capacity fixes \(H_{\mathrm{dS}}\), and the flat Friedmann constraint at
\(a=1\) gives
\begin{equation}
  \Omega_{A0}
  =
  1-\Omega_{\Lambda,\OPH}-\Omega_{b0}-\Omega_{\nu0}-\Omega_{r0},
  \qquad
  \Omega_{\Lambda,\OPH}
  =
  \left(\frac{H_{\mathrm{dS}}}{H_0}\right)^2.
  \label{eq:flatresidual}
\end{equation}
The transported homogeneous dust branch is
\begin{equation}
  \rho_A(a)=\rho_{A0}a^{-3}.
\end{equation}
Using \(H_0=67.4\,\mathrm{km\,s^{-1}\,Mpc^{-1}}\),
\(\Omega_b h^2=0.02237\), the OPH neutrino mass sum, and
\(\Omega_r=9.17\times10^{-5}\), the residual is:
\begin{center}
\begin{tabular}{@{}lr@{}}
\toprule
Quantity & Value \\
\midrule
\(H_{\mathrm{dS}}\) & \(55.759940256\,\mathrm{km\,s^{-1}\,Mpc^{-1}}\) \\
\(\Omega_{\Lambda,\OPH}\) & \(0.684423332\) \\
\(\Omega_b\) & \(0.049243191\) \\
\(\Omega_\nu\) & \(0.002127375\) \\
\(\Omega_r\) & \(0.000091700\) \\
\(\Omega_A\) & \(0.264114401\) \\
\(\Omega_m\) & \(0.315484968\) \\
\(\rho_{A0}\) & \(2.253649933\times10^{-27}\,\mathrm{kg\,m^{-3}}\) \\
\bottomrule
\end{tabular}
\end{center}
This matches the Planck-like cold-matter budget at the compressed-parameter
level.  The number is a state-selection residual.  The static galaxy RAR law
does not determine it.  The OPH theorem that would make it standalone is a
finite-screen additive-load selector for the conserved homogeneous anomaly
charge.

The closed transported branch does give a theorem-grade current and charge.
Let \(u_A^a\) be the branch four-velocity and let \(\rho_A\) be the
coarse-grained information-defect load density.  The load current is
\begin{equation}
  J_A^a=\rho_A u_A^a .
\end{equation}
On the closed branch,
\begin{equation}
  \nabla_aJ_A^a=0,
  \qquad
  Q_A[\Sigma]=\int_{\Sigma} J_A^a n_a\,\dd\Sigma
\end{equation}
is independent of the homogeneous slice.  In FLRW this gives
\begin{equation}
  Q_A=a^3\rho_A V_{\mathrm{com}},
  \qquad
  \rho_A(a)=\frac{Q_A}{a^3V_{\mathrm{com}}}.
\end{equation}
This proves transport of a supplied homogeneous load.  It does not fix the
amplitude \(Q_A\).  The missing OPH object is a nontrivial finite-screen
additive load invariant
\begin{equation}
  Q_A=
  F_{\OPH}\!\left(
    N_{\mathrm{scr}},P,\text{screen normal-form data}
  \right).
\end{equation}
Without that invariant, the abundance is cosmological state data or the flat
capacity-saturated residual \eqref{eq:flatresidual}.

\subsection{Environmental Kernel}

Around a nonzero local background field
\(\bar{\mathbf g}_b\), write
\begin{equation}
  x=\frac{|\bar{\mathbf g}_b|}{\aoph},
  \qquad
  \mathbf n=\frac{\bar{\mathbf g}_b}{|\bar{\mathbf g}_b|}.
\end{equation}
Linearizing the anomaly field
\((\nu_{\OPH}-1)\mathbf g_b\) gives the susceptibility
\begin{equation}
  A_{ij}
  =
  [\nu_{\OPH}(x)-1]\delta_{ij}
  +x\nu_{\OPH}'(x)n_i n_j,
\end{equation}
where
\begin{equation}
  \nu_{\OPH}'(x)
  =
  -\frac{
    \lambdac\exp[-\lambdac\sqrt{x}]
  }{
    2\sqrt{x}[1-\exp(-\lambdac\sqrt{x})]^2
  }.
\end{equation}
For a Fourier mode whose direction has cosine \(\mu\) against \(\mathbf n\),
the density-response multiplier is
\begin{equation}
  K_A^{(\rho)}(x,\mu)
  =
  \nu_{\OPH}(x)-1+x\nu_{\OPH}'(x)\mu^2.
  \label{eq:environmentalkernel}
\end{equation}
An isotropic environmental average gives
\begin{equation}
  \langle K_A^{(\rho)}\rangle
  =
  \nu_{\OPH}(x)-1+\frac{x}{3}\nu_{\OPH}'(x).
\end{equation}
The contrast kernel used in perturbation equations is normalized by the
background density ratio:
\begin{equation}
  B_A(k,a)
  =
  \frac{\bar\rho_b(a)}{\bar\rho_A(a)}K_A^{(\rho)}(k,a).
\end{equation}
In a realistic cosmological calculation, this local expression has to be
integrated against the distribution of environmental fields,
\begin{equation}
  K_A^{(\rho)}(k,a)
  =
  \int\dd x\,\dd\mu\,
  \Pi(x,\mu|k,a)
  [\nu_{\OPH}(x)-1+x\nu_{\OPH}'(x)\mu^2].
\end{equation}
This average is finite only under an explicit small-field support condition,
for example
\begin{equation}
  \int_0^\epsilon \dd x\,\dd\mu\,
  \Pi(x,\mu|k,a)\,x^{-1/2}<\infty
  \qquad
  \text{for some }\epsilon>0.
\end{equation}
If \(\Pi\) has an exact atom at \(x=0\), the local nonzero-field kernel is not
the correct FLRW response.
The distribution \(\Pi\) is an OPH finite-collar output or an explicitly
declared environmental closure.  Exact FLRW has \(x=0\), so
\eqref{eq:environmentalkernel} does not provide a regular universal CMB
density-response kernel by itself.

The implementation contract uses one parent finite-collar functional:
\begin{equation}
  \rho_{A,\mathrm{eq}}[X]c^2
  =
  \frac{15}{8\pi^2\ell(X)^4}
  \int_{\mathcal C_X}\dd\mu_C\,
  I_{\omega_C}(A:D|B).
  \label{eq:parentfunctional}
\end{equation}
For finite samples \(s\) with weights \(w_s\), define
\begin{align}
  R(X)&=\frac{\sum_s w_s I_s}{\sum_s w_s},\\
  K_A^{(\rho)}(k,a)
  &=\frac{1}{\bar\rho_b}
    \frac{\partial \rho_{A,\mathrm{eq}}}{\partial \delta_b(k,a)},\\
  B_A(k,a)
  &=\frac{\bar\rho_b(a)}{\bar\rho_A(a)}K_A^{(\rho)}(k,a).
\end{align}
The same evaluator emits the homogeneous equilibrium density, the linear
density-response kernel, the settled galaxy source, and the cluster
equilibrium source.  A theorem-grade prediction requires the OPH collar
measure \(\dd\mu_C\), the finite-screen ensemble, and the small-field support
condition.  Fitting \(\Pi\) to CMB, weak-lensing, SPARC, or cluster data is
excluded.

\subsection{Compressed CMB, BAO, Growth, And \(S_8\) Diagnostic}

The local Boltzmann plumbing test passes the anomaly as a cold pressureless
component with the flat-residual diagnostic parent-grid ratio
\(\rho_A/\rho_b=5.363470441\).  It uses CAMB \(1.6.6\), the OPH neutrino mass
sum, and diagonal Gaussian compressed rows:
\begin{center}
\begin{tabular}{@{}lrrr@{}}
\toprule
Quantity & Prediction & Target & Pull \\
\midrule
Planck \(\Omega_m\) & \(0.315905207\) & \(0.315\pm0.007\) & \(+0.129\) \\
Planck \(\sigma_8\) & \(0.807787208\) & \(0.811\pm0.006\) & \(-0.535\) \\
Planck \(S_8\) & \(0.828924043\) & \(0.831027707\pm0.0111\) & \(-0.190\) \\
DESI DR1 BAO \(\Omega_m\) & \(0.315905207\) & \(0.295\pm0.015\) & \(+1.394\) \\
DESI DR1 BAO/BBN/\(\theta_\ast\) \(H_0\) & \(67.4\) & \(68.52\pm0.62\) & \(-1.806\) \\
Weak-lensing \(S_8\) & \(0.828924043\) & \(0.790\pm0.016\) & \(+2.433\) \\
\bottomrule
\end{tabular}
\end{center}
The diagonal compressed value is \(\chi^2=11.463\) for six rows.  This is a
plumbing check.  It does not replace Planck, DESI, or weak-lensing likelihoods.
A publication-grade calculation needs the OPH \(B_A(k,a)\) kernel in a custom
Boltzmann module and the full experimental covariances.

\section{Likelihood Contracts}

\subsection{SPARC Disk-Potential Contract}

The publication-grade SPARC test must solve the conservative disk equation
\begin{equation}
  \nabla^2\Phi
  =
  \nabla\cdot
  \left[
    \nu_\lambda\!\left(\frac{|\nabla\Phi_b|}{\aoph}\right)
    \nabla\Phi_b
  \right]
\end{equation}
on an axisymmetric baryonic disk model.  The hierarchy must marginalize or
profile distance, inclination, disk and bulge mass-to-light ratios, gas scale,
intrinsic scatter, velocity covariance, quality cuts, and grid convergence.
The unit branch and the \(\Zsix\)/Poisson branch are fixed before this test;
SPARC is used as a comparison, not as the selector for \(\lambdac\).

\subsection{Cluster Map Contract}

The cluster forward model uses gas hydrodynamics, stellar and galaxy tracers,
transported defect stress, the repair source
\(\rho_{A,\mathrm{eq}}\), lensing projection, X-ray emissivity, SZ pressure,
and covariance-aware map likelihoods.  Merging clusters test whether the
transported stress lags collisional gas.  Relaxed clusters test whether the
same stress law returns to the static equilibrium source.  The static-only
failure mode is part of the likelihood contract.

\subsection{Boltzmann Contract}

The Boltzmann module exposes
\[
  \bar\rho_A(a),\quad
  \bar\rho_{A,\mathrm{eq}}(a),\quad
  w_A(a),\quad
  c_{s,A}^2(k,a),\quad
  \sigma_A(k,a),\quad
  Q_A^\mu,\quad
  B_A(k,a),\quad
  \Gamma_{\mathrm{rec}}(k,a),
\]
with the OPH neutrino mass sum in the same run.  Every entry must be derived
from a frozen finite covariant parent artifact rather than fitted from CMB,
BAO, lensing, SPARC, or cluster data.  It must recover the \(\LCDM\)
cold-component limit when the exchange and stress corrections are turned off,
reproduce the compressed diagnostic rows only as a plumbing check, and then run
the full CMB TT/TE/EE, CMB lensing, BAO, supernova, weak-lensing, RSD, and
\(S_8\) likelihoods under the same nuisance model used for \(\LCDM\) and
\(w_0w_a\).  The physical likelihood run is frozen by immutable hashes for the
source artifact, Boltzmann solver, and likelihood configuration before the data
comparison.

\section{Numerical Galaxy Comparison}

The OPH capacity benchmark used here is
\begin{equation}
  \aoph = 1.029186271\times10^{-10}\,\mathrm{m\,s^{-2}}.
\end{equation}
The common empirical reference is
\begin{equation}
  a_0^{\mathrm{emp}}=1.2\times10^{-10}\,\mathrm{m\,s^{-2}}.
\end{equation}
The \(\Zsix\)/Poisson branch gives
\begin{equation}
  \lambdac=0.934300639,
  \qquad
  \aeff=1.179018696\times10^{-10}\,\mathrm{m\,s^{-2}}.
\end{equation}

\begin{center}
\begin{tabular}{@{}lrr@{}}
\toprule
Quantity & Value & Comparison \\
\midrule
\(\aoph\) &
\(1.029186271\times10^{-10}\) &
\(-8.070\%\) versus best same-function RAR \\
\(\aoph\) &
\(1.029186271\times10^{-10}\) &
\(-14.234\%\) versus empirical reference \\
\(\aeff\), \(\Zsix\)/Poisson &
\(1.179018696\times10^{-10}\) &
\(+5.314\%\) versus best same-function RAR \\
\(\aeff\), \(\Zsix\)/Poisson &
\(1.179018696\times10^{-10}\) &
\(-1.748\%\) versus empirical reference \\
\bottomrule
\end{tabular}
\end{center}

\subsection{SPARC RAR}

Using the public SPARC RAR data \cite{sparcpage,mcgaugh2016}, the local
diagnostic gives:
\begin{center}
\begin{tabular}{@{}lrrr@{}}
\toprule
Scenario & \(a_0\) in \(\mathrm{m\,s^{-2}}\) & RMS dex & Binned RMS dex \\
\midrule
OPH unit & \(1.029186271\times10^{-10}\) & \(0.134208\) & \(0.026169\) \\
\(\Zsix\)/Poisson & \(1.179018696\times10^{-10}\) & \(0.132834\) & \(0.015760\) \\
Empirical reference & \(1.200000000\times10^{-10}\) & \(0.132913\) & \(0.016164\) \\
Best same-function & \(1.119530397\times10^{-10}\) & \(0.132947\) & \(0.017428\) \\
\bottomrule
\end{tabular}
\end{center}

\subsection{Rotation Diagnostics}

The fixed stellar mass-to-light diagnostic and the first nuisance-profiled
diagnostic give:
\begin{center}
\begin{tabular}{@{}lrrr@{}}
\toprule
Scenario & Fixed-M/L RMS & Profiled-M/L \(\chi^2/\mathrm{pt}\) & Profiled RMS \\
\midrule
OPH unit & \(23.366208\) & \(1.341948\) & \(12.420181\) \\
\(\Zsix\)/Poisson & \(22.868866\) & \(1.390941\) & \(12.590203\) \\
Empirical reference & \(22.842145\) & \(1.400930\) & \(12.624079\) \\
\bottomrule
\end{tabular}
\end{center}
Units for the RMS columns are \(\mathrm{km\,s^{-1}}\).

\subsection{Systematic Likelihood Scaffold}

A stronger local scaffold profiles stellar M/L, distance, inclination, and gas
scale per galaxy using the SPARC sample table \cite{lelli2016,sparcpage}.  It
uses \(Q\le2\), inclination at least \(30^\circ\), and 3168 rotation-curve
rows.
\begin{center}
\begin{tabular}{@{}lrrrr@{}}
\toprule
Scenario & \(\chi^2/\mathrm{pt}\) & \(\chi^2/\mathrm{dof}\) & RMS & Median disk M/L \\
\midrule
OPH unit & \(1.016589\) & \(1.275672\) & \(11.249541\) & \(0.519102\) \\
\(\Zsix\)/Poisson & \(1.021093\) & \(1.281324\) & \(11.276832\) & \(0.497435\) \\
Empirical reference & \(1.023261\) & \(1.284045\) & \(11.285518\) & \(0.493865\) \\
\bottomrule
\end{tabular}
\end{center}
The systematic scaffold mildly favors the unit branch.  The \(\Zsix\)/Poisson
branch remains the better acceleration-scale and fixed-M/L RAR match.  The
fixed-M/L improvement alone is weak evidence for the \(\Zsix\) bridge theorem.

\subsection{Empirical Implementation Scorecard}

The executable empirical pass combines four diagnostic layers: public SPARC
tables, the flat capacity-saturated homogeneous anomaly state, the finite
repair matrix with a cluster timing gate, and the compressed CAMB rows.  The
single-pass output is:
\begin{center}
\small
\begin{tabularx}{\textwidth}{@{}p{0.39\textwidth}p{0.24\textwidth}X@{}}
\toprule
Quantity & OPH output & Measurement comparison \\
\midrule
\(\Omega_A\) & \(0.264114401\) & Planck-like matter budget \\
\(\Omega_{\Lambda,\OPH}\) & \(0.684423332\) & capacity de Sitter branch \\
\(\rho_A/\rho_b\) & \(5.363470441\) & flat residual ratio \\
\(\gamma_{\mathrm{rec}}\) & \(0.050717471\) & finite repair matrix gap \\
\(\aoph\) & \(1.029186271\times10^{-10}\) & \(8.070\%\) low versus best RAR \\
\(\aeff\), \(\Zsix\)/Poisson & \(1.179018696\times10^{-10}\) & \(1.748\%\) low versus common empirical \(a_0\) \\
SPARC systematic \(\chi^2/\mathrm{pt}\), unit & \(1.016589\) & best scaffold row \\
SPARC systematic \(\chi^2/\mathrm{pt}\), \(\Zsix\) & \(1.021093\) & close scaffold row \\
CAMB \(\Omega_m\) & \(0.315905206\) & \(+0.129\sigma\) versus Planck \\
CAMB \(\sigma_8\) & \(0.807787204\) & \(-0.535\sigma\) versus Planck \\
CAMB \(S_8\) & \(0.828924037\) & \(+2.433\sigma\) versus weak lensing \\
\bottomrule
\end{tabularx}
\end{center}

The compressed CMB/BAO/growth score is
\[
  \chi^2_{\mathrm{diag}}=11.463
  \qquad \hbox{for six compressed rows.}
\]
The pressure point is weak-lensing \(S_8\).  The diagnostic value is
\(0.828924037\).  The DES/KiDS compressed row is \(0.790\pm0.016\).

For the cluster timing gate, take a \(200\,\mathrm{kpc}\) initial anomaly
separation, a \(0.2\,\mathrm{Gyr}\) time after passage, and an example measured
offset \(150\pm50\,\mathrm{kpc}\).  The repair matrix gives:
\begin{center}
\begin{tabular}{@{}rrrrr@{}}
\toprule
\(\tau_{\mathrm{rec}}\) in Gyr & \(t_{\mathrm{commit}}\) in Myr &
Retained fraction & Offset in kpc & \(\chi^2\) \\
\midrule
\(0.3\) & \(15.2152\) & \(0.513417\) & \(102.683\) & \(0.896\) \\
\(0.5\) & \(25.3587\) & \(0.670320\) & \(134.064\) & \(0.102\) \\
\(1.0\) & \(50.7175\) & \(0.818731\) & \(163.746\) & \(0.076\) \\
\bottomrule
\end{tabular}
\end{center}
This row is a timing gate.  Map likelihoods need observed lensing maps, gas
maps, stellar maps, merger-age priors, and covariance.

\subsection{Implementation Work Packages}

The code required for paper-grade empirical tests consists of the following
modules:
\begin{enumerate}
\item An axisymmetric conservative disk-potential solver for
\[
  \nabla^2\Phi
  =
  \nabla\cdot[\nu(|\nabla\Phi_b|/\aoph)\nabla\Phi_b],
\]
with grid-convergence tests and lensing-potential output.
\item A hierarchical SPARC likelihood that samples or profiles distance,
inclination, disk M/L, bulge M/L, gas scale, intrinsic scatter, velocity
covariance, and selection cuts under one nuisance model.
\item A finite-collar parent evaluator for \(\rho_A(a)\),
\(K_A^{(\rho)}(k,a)\), \(B_A(k,a)\), and
\(\rho_{A,\mathrm{eq}}[X]\), using OPH collar samples as input.  The
scalar-load diagnostic grid is plumbing only.
\item A CAMB or CLASS anomaly module exposing the background density, stress
variables, exchange current, relaxation kernel, environmental response
kernel, and the OPH neutrino mass sum in one run.
\item A cluster forward model with gas hydrodynamics, galaxy and stellar
tracers, transported anomaly stress, lensing projection, X-ray and SZ
observables, merger-age priors, and map covariance.
\item A likelihood runner for Planck or ACT spectra, CMB lensing, BAO,
supernovae, weak lensing, RSD, SPARC, and cluster maps, with model comparison
against \(\LCDM\), empirical MOND functions, halo models, and \(w_0w_a\)
extensions.
\item A reproducibility harness with fixed data hashes, fixed random seeds,
unit checks, and independent numerical replication of the scorecard rows.
\end{enumerate}

\subsection{Correction Audit}

The comparison mostly tests normalization because the OPH static law uses the
same acceleration-family shape as the empirical RAR fit.  If \(\aoph\) is held
fixed, different targets imply different \(\lambdac\) values:

\begin{center}
\begin{tabular}{@{}lrrr@{}}
\toprule
Target & Required \(\lambdac\) & Reserve \(-\ln\lambdac\) & Ratio to \(P/24\) \\
\midrule
Best same-function RAR & \(0.958802256\) & \(0.042070424\) & \(0.619074\) \\
Binned RAR & \(0.936737944\) & \(0.065351711\) & \(0.961663\) \\
Common empirical \(a_0\) & \(0.926096769\) & \(0.076776547\) & \(1.129781\) \\
Fixed-M/L rotation & \(0.926654674\) & \(0.076174302\) & \(1.120919\) \\
OPH \(\Zsix\)/Poisson & \(0.934300639\) & \(0.067957009\) & \(1.000000\) \\
\bottomrule
\end{tabular}
\end{center}

This rules out several tempting shortcuts.  In the normalization-profiled
all-point objective, a generalized activation exponent near \(0.503420\)
appears.  With fixed \(\aoph\), direct RAR objectives give exponents near
\(0.516\) all-point and \(0.519\) weighted-bin.  These shifts are
objective-dependent, and flat deep-IR rotation requires \(0.5\).  Moving the
exponent would buy little and would damage the BTFR theorem.

Finite-thickness nonuniformity can raise \(\lambdac\) above \(\exp(-P/24)\),
which moves the coefficient toward the binned RAR value.  It cannot lower
\(\lambdac\) toward the common empirical or fixed-M/L values while the mean
reserve remains \(P/24\) and Poisson counting is retained.  An exact hit to
the common empirical acceleration would require about \(13\%\) more protected
inactive reserve than \(P/24\), or a different microscopic selector.  That
extra reserve needs an OPH theorem before use.

\section{Claim Boundaries}

\begin{longtable}{@{}p{0.30\textwidth}p{0.46\textwidth}p{0.16\textwidth}@{}}
\toprule
Pressure point & OPH route & Status \\
\midrule
\endhead
Unknown particle & Dark source is a carried modular/collar information-defect
stress. & Conditional source target. \\
Galaxy RAR and BTFR & Collar/cut activation gives the square-root transition
and Poisson activation law under independent-increment repair counting. &
Conditional static theorem. \\
Acceleration scale & Capacity scale gives \(\aoph\); \(\Zsix\)/Poisson reserve
gives \(\aeff\) close to the empirical scale. & Conditional coefficient
theorem. \\
Galaxy lensing & No-slip branch gives \(\Phi=\Psi\) if the information-defect
stress is dust on settled galaxies. & Conditional stress theorem. \\
Cluster offsets & Transported anomaly can lag collisional gas; the Jacobi
clock gives a no-fit scale under the tidal-mismatch premise. & Forward
likelihood contract. \\
CMB and growth & Flat capacity state selection gives a CDM-like homogeneous
residual; the closed branch has a conserved load current; environmental
kernels describe nonzero-field response. & Amplitude selector and real
Boltzmann likelihood work in progress. \\
\bottomrule
\end{longtable}

The supported claim is:
\begin{quote}
OPH defines an imperfect-information correction to the recovered Einstein
branch.  Under explicit bridge premises, the static galaxy equation reproduces
the RAR/BTFR functional form from collar/cut repair statistics and predicts an
acceleration scale close to SPARC.  The unit branch is low in normalization.
The \(\Zsix\)/Poisson branch lands within \(1.748\%\) of the common empirical
acceleration reference.  Profiled stellar, distance, inclination, and gas
nuisances mildly favor the unit branch.  The Jacobi clock gives a no-fit
cluster-offset scale under its tidal-mismatch premise.  Flat capacity state
selection gives a Planck-like homogeneous information-defect residual, while
closed-branch transport gives a conserved load current.  Cluster and
CMB/growth likelihoods require measured cluster covariances, a finite-screen
amplitude selector, a finite covariant source parent, recipient stress for any
nonzero repair exchange, and a frozen Boltzmann/likelihood implementation.
\end{quote}

Claim outside this paper:
\begin{quote}
OPH has fully derived dark matter and solved MOND's cluster and CMB problems.
\end{quote}
That stronger statement requires SPARC publication likelihoods, cluster
likelihoods, a finite covariant parent for the cosmological source functions,
and frozen Boltzmann likelihoods.

\section{Proof Ledger}

\begin{center}
\begin{tabularx}{\textwidth}{@{}p{0.25\textwidth}Xp{0.22\textwidth}@{}}
\toprule
Target & OPH closure & Numerical state \\
\midrule
Galaxy source & Minimal scalar recoverability defect
\(R_C=I(A:D|B)\). & Conditional source target. \\
Quotient-edge locality & Scalar repair opportunities are edge-center
cut-register events on the physical quotient. & Conditional bridge. \\
Finite thickness & Product transverse regulator with one quotient trace makes
the protected center reserve uniform through the weighted collar. &
Conditional bridge. \\
No-slip stress & Positive kinetic collar measure; settled monokinetic branch
has dust stress and vanishing information-sector anisotropic stress. &
Conditional stress theorem. \\
Repair relaxation & Fixed repair-commit chain gives
\(\tau_{\mathrm{rec}}=t_{\mathrm{commit}}/\gamma_{\mathrm{rec}}\). &
Conditional packet theorem and clock candidate. \\
CMB and growth & Nonzero-environment kernel is explicit; exact FLRW static
RAR kernel is blocked by a no-go theorem. &
Finite-parent Boltzmann contract. \\
Abundance history & Flat capacity state selection gives a CDM-like residual;
the closed transported branch has a conserved load current and charge. &
Amplitude selector work in progress. \\
Linear transfer & The perturbation equation has a relaxation transfer
solution between CDM-like and tracking limits. & Requires evaluated
\(B_A(k,a)\), recipient stress for nonzero exchange, and frozen hashes. \\
Parent functional & The finite collar expectation of the positive repair
defect emits \(\rho_{A,\mathrm{eq}}\) and \(B_A\) if its FLRW limit is
regular. & Scalar source target, not full stress closure. \\
Repair matrix & A reversible Metropolis kernel on packet states gives a
concrete \(K\). & Conditional on declared packet schema. \\
Jacobi clock & The tidal scalar \((E_{ij}E^{ij}/6)^{1/4}\) gives a no-fit
dynamic repair scale under the Jacobi mismatch premise. & Conditional clock. \\
\bottomrule
\end{tabularx}
\end{center}

\subsection{Positive Galaxy Source}

\paragraph{Theorem target 0A.}
Let a settled galaxy collar state be judged by a scalar repair defect \(D_C\)
with these properties: \(D_C\ge0\); \(D_C=0\) exactly on locally recoverable
Markov states; \(D_C\) is invariant under boundary redundancy; \(D_C\) is
additive over independent collar components; and \(D_C\) is monotone under
local recovery maps.  Then the minimal scalar branch is
\begin{equation}
  D_C=\kappa I(A:D|B),
  \qquad
  \kappa\ge0.
\end{equation}
The OPH dark branch uses \(R_C=I(A:D|B)\).

\paragraph{Explicit premise.}
The galaxy-scale finite-stage modular nonadditivity is represented by this
minimal scalar recoverability defect on settled collar states.

\paragraph{Conditional argument.}
Conditional mutual information has all listed properties on a finite collar.
A different scalar with the same zero set and additivity would add an
independent scalar datum to the minimal branch.  The minimal scalar selector is
therefore fixed up to a nonnegative multiplier.  The OPH unit convention sets
that multiplier to one.  Since \(I(A:D|B)\ge0\), the weak-field density defined
by the anomaly normalization is nonnegative. \(\square\)

\paragraph{Proof obligation.}
Derive the map from the finite-stage modular nonadditivity term to
\(I(A:D|B)\) on the galaxy branch.

\paragraph{Failure consequence.}
If the modular remainder is not this positive defect, the attractive source and
normalization formula must be replaced by the derived scalar.

\subsection{Quotient-Edge Scalar Opportunity}

\paragraph{Theorem target 0B.}
At fixed cutoff, physical collar observables live on the quotient-local
fixed-point algebra.  A settled galaxy scalar repair opportunity is the local
event that the quotient-local recoverability defect on a collar cut is
nonzero.  Then the scalar opportunity is counted on the edge-center cut
register.

\paragraph{Explicit premise.}
The scalar recoverability event has no independent bulk, screen-area, or
noncentral collar channel.

\paragraph{Conditional argument.}
The defect is a physical scalar.  Boundary-redundant representatives produce
the same defect value because the quotient repair law descends through the
projection onto physical collar states.  The event that the defect is nonzero
is therefore an observable event on the quotient-local collar algebra.  For a
connected cut, the fixed-cutoff edge-center decomposition supplies the central
projectors that resolve physical cut sectors.  A scalar event with no vector,
tensor, or bulk label selects support through those central projectors.  Thus
the scalar opportunity count is a quotient-edge cut-register count.  The
protected \(\Zsix\) center reserve is supported on the same register.
\(\square\)

\paragraph{Proof obligation.}
Prove the fixed-cutoff uniqueness statement: every quotient-local scalar
recoverability event is resolved on the edge-center cut register.

\paragraph{Failure consequence.}
If a separate scalar channel exists, the \(\Zsix\) reserve does not fix
\(\lambdac\) without an additional co-registration theorem.

\subsection{Uniform Finite-Thickness Reserve}

\paragraph{Theorem target 0C.}
Let the finite collar be a transverse thickening of one quotient-edge cut
algebra with normalized transverse weight \(w(y)\) and a single quotient trace
on the edge algebra.  Let the protected \(\Zsix\) reserve be a central edge
projector whose mean on that trace is \(\epsilon_{\Zsix}=P/24\).  Then
\begin{equation}
  \lambdac=\exp(-P/24).
\end{equation}

\paragraph{Explicit premise.}
The scalar-weighted transverse regulator is a product thickening with one
quotient trace and uniform protected-center density.

\paragraph{Conditional argument.}
The protected reserve belongs to the edge-center algebra.  The transverse
coordinate belongs to the regulator thickening.  With one quotient trace on the
edge algebra, the trace of the protected central projector is the same on every
transverse slice that carries scalar activation weight.  Hence
\(\epsilon_{\Zsix}(y)=\epsilon_{\Zsix}\).  Poisson zero-reserve survival gives
\(\exp(-\epsilon_{\Zsix})\) at each slice.  Averaging with \(w(y)\) leaves the
same value.  Using \(\epsilon_{\Zsix}=P/24\) gives the stated coefficient.
\(\square\)

\paragraph{Proof obligation.}
Derive \(w(y)\), the transverse regulator, and the quotient trace from base OPH
screen/collar data.

\paragraph{Failure consequence.}
If the reserve density is nonuniform, the coefficient is
\[
  \int\dd y\,w(y)\exp[-\epsilon_{\Zsix}(y)],
\]
not necessarily \(\exp(-P/24)\).

\subsection{No-Slip Settled Stress}

\paragraph{Theorem target 0D.}
Let the information-defect stress be the second moment of a positive
finite-collar kinetic measure \(\dd\mathcal E_x(v)\) on the future unit
timelike shell:
\begin{equation}
  T_{\mathcal I}^{ab}=\int v^a v^b\,\dd\mathcal E_x(v).
\end{equation}
If the settled galaxy branch is monokinetic in its rest frame, then
\begin{equation}
  T_{\mathcal I}^{ab}
  =
  \rho_{\mathcal I}u_{\mathcal I}^a u_{\mathcal I}^b,
  \qquad
  P_{\mathcal I}=0,
  \qquad
  q_{\mathcal I}^a=0,
  \qquad
  \pi_{\mathcal I}^{ab}=0.
\end{equation}
For negligible or separately modeled anisotropic stress in all other sectors,
\(\Phi=\Psi\).

\paragraph{Explicit premise.}
The finite collar packet state emits the positive kinetic measure, and the
settled galaxy branch is monokinetic at the coarse scale used for lensing.

\paragraph{Conditional argument.}
A positive measure on future timelike velocities has a symmetric
future-positive second moment.  In the monokinetic branch the measure is
concentrated at one four-velocity, so the second moment is
\(\rho_{\mathcal I}u_{\mathcal I}^a u_{\mathcal I}^b\).  Its pressure, heat
flux, and anisotropic stress vanish in that rest frame.  The off-diagonal
spatial Einstein equation then gives \(\Phi=\Psi\) when the other sectors
carry no uncompensated anisotropic stress. \(\square\)

\paragraph{Proof obligation.}
Derive \(\dd\mathcal E_x(v)\) from the finite OPH collar packet state or from
an equivalent covariant parent action.

\paragraph{Failure consequence.}
If pressure or anisotropic stress survives coarse graining, the lensing slip
equation must include those terms and no-slip is regime-dependent.

\subsection{Repair Relaxation Scale}

\paragraph{Theorem 0E.}
Let \(K\) be the fixed-cutoff repair-commit transition matrix for one local
coarse cell at fixed baryonic source.  If \(K\) is finite, irreducible, and
aperiodic, with spectral gap \(\gamma_{\mathrm{rec}}\), then
\begin{equation}
  \tau_{\mathrm{rec}}
  =\frac{t_{\mathrm{commit}}}{\gamma_{\mathrm{rec}}}.
\end{equation}
If every allowed repair move has probability at least \(q_{\min}>0\) and the
comparison chain has conductance \(\Phi_K\), then
\begin{equation}
  \gamma_{\mathrm{rec}}\ge \frac12 q_{\min}\Phi_K^2.
\end{equation}

\paragraph{Proof.}
A finite irreducible aperiodic chain has a unique stationary state.  Every
nonstationary mode decays as \(\exp(-\gamma_n n)\) in repair-step time, and
the slowest nonzero mode has rate \(\gamma_{\mathrm{rec}}\).  One physical
commit step takes \(t_{\mathrm{commit}}\), so physical relaxation has rate
\(\gamma_{\mathrm{rec}}/t_{\mathrm{commit}}\).  The conductance bound follows
from the standard Cheeger comparison for the reversible comparison chain and
the lower move probability \(q_{\min}\). \(\square\)

\subsection{Environmental Kernel Around A Nonzero Field}

\paragraph{Conditional theorem 0F.}
Let the settled anomaly field be
\begin{equation}
  \mathbf g_A=(\nu_{\OPH}-1)\mathbf g_b.
\end{equation}
Linearize around a nonzero background field
\(\bar{\mathbf g}_b\).  With
\[
  x=|\bar{\mathbf g}_b|/\aoph,
  \qquad
  n_i=\bar g_{b,i}/|\bar{\mathbf g}_b|,
\]
the Fourier density response to a baryonic density mode is
\begin{equation}
  K_A^{(\rho)}(x,\mu)
  =
  \nu_{\OPH}(x)-1+x\nu_{\OPH}'(x)\mu^2,
\end{equation}
where \(\mu\) is the cosine between the mode direction and \(\mathbf n\).
The contrast kernel is
\begin{equation}
  B_A(k,a)
  =
  \frac{\bar\rho_b(a)}{\bar\rho_A(a)}K_A^{(\rho)}(k,a).
\end{equation}

\paragraph{Proof.}
The first variation is
\begin{equation}
  \delta g_{A,i}
  =
  \left(
    [\nu_{\OPH}(x)-1]\delta_{ij}
    +x\nu_{\OPH}'(x)n_i n_j
  \right)\delta g_{b,j}.
\end{equation}
Taking the divergence and using the Fourier-space Poisson relation for
\(\delta\mathbf g_b\) gives the displayed scalar multiplier.  The derivative
of
\[
  \nu_{\OPH}(x)=[1-\exp(-\lambdac\sqrt{x})]^{-1}
\]
is
\[
  \nu_{\OPH}'(x)
  =
  -\frac{\lambdac\exp[-\lambdac\sqrt{x}]}
  {2\sqrt{x}[1-\exp(-\lambdac\sqrt{x})]^2}.
\]
\(\square\)

\subsection{Background Abundance Green Function}

\paragraph{Theorem 0G.}
Let the background anomaly density obey
\begin{equation}
  \rho_A' + 3\mathcal H\rho_A
  = -a\Gamma_{\mathrm{rec}}(\rho_A-\rho_{A,\mathrm{eq}}).
\end{equation}
Define the comoving densities
\begin{equation}
  R_A=a^3\rho_A,
  \qquad
  R_{A,\mathrm{eq}}=a^3\rho_{A,\mathrm{eq}}.
\end{equation}
Then
\begin{equation}
  R_A(\eta)=
  W(\eta,\eta_i)R_A(\eta_i)
  +\int_{\eta_i}^{\eta}\dd\eta'\,
    W(\eta,\eta')a(\eta')\Gamma_{\mathrm{rec}}(\eta')
    R_{A,\mathrm{eq}}(\eta'),
\end{equation}
where
\begin{equation}
  W(\eta,\eta')
  =
  \exp\!\left[
    -\int_{\eta'}^{\eta}\dd\tilde\eta\,
      a(\tilde\eta)\Gamma_{\mathrm{rec}}(\tilde\eta)
  \right].
\end{equation}

\paragraph{Proof.}
Multiplying the background equation by \(a^3\) gives
\begin{equation}
  R_A'
  =
  -a\Gamma_{\mathrm{rec}}(R_A-R_{A,\mathrm{eq}}).
\end{equation}
This is a first-order linear equation.  Multiplication by the integrating
factor
\begin{equation}
  \exp\!\left[
    \int_{\eta_i}^{\eta}\dd\eta'\,
      a(\eta')\Gamma_{\mathrm{rec}}(\eta')
  \right]
\end{equation}
and integration gives the stated expression. \(\square\)

The theorem shows exactly where abundance information enters: the initial
comoving anomaly load, the equilibrium source, and the repair rate.  These are
OPH microphysical outputs, not SPARC parameters.

\subsection{Linear Relaxation Transfer}

\paragraph{Theorem 0H.}
For fixed \(k\), assume the no-slip anomaly perturbation equation
\begin{equation}
  \delta_A'
    =-\theta_A+3\Phi'
      -a\Gamma_{\mathrm{rec}}q_A
       (\delta_A-B_A\delta_b).
\end{equation}
Define
\begin{equation}
  \Xi_A=a\Gamma_{\mathrm{rec}}q_A.
\end{equation}
Then
\begin{equation}
\begin{aligned}
  \delta_A(\eta,k)
  &=W_A(\eta,\eta_i)\delta_A(\eta_i,k)\\
  &\quad
  +\int_{\eta_i}^{\eta}\dd\eta'\,
    W_A(\eta,\eta')
    \left[
      -\theta_A+3\Phi'
      +\Xi_A B_A\delta_b
    \right]_{\eta',k},
\end{aligned}
\end{equation}
with
\begin{equation}
  W_A(\eta,\eta')
  =
  \exp\!\left[
    -\int_{\eta'}^{\eta}\dd\tilde\eta\,\Xi_A(\tilde\eta,k)
  \right].
\end{equation}

\paragraph{Proof.}
Move the damping term to the left:
\begin{equation}
  \delta_A' + \Xi_A\delta_A
  =
  -\theta_A+3\Phi' + \Xi_A B_A\delta_b.
\end{equation}
The integrating-factor solution is the displayed expression. \(\square\)

This formula makes the two useful limits explicit.  If
\(\Xi_A/\mathcal H\ll1\), the anomaly evolves like pressureless free-falling
matter.  If \(\Xi_A/\mathcal H\gg1\), it tracks the OPH equilibrium source up
to derivative corrections.  Cluster offsets and settled galaxies require the
same relaxation operator to sit between those limits.

\subsection{Static-To-Linear Separation}

\paragraph{Theorem 0I.}
The nonlinear static RAR law
\begin{equation}
  g_{\mathrm{obs}}
  =
  \frac{g_b}{1-\exp[-\lambdac\sqrt{g_b/\aoph}]}
\end{equation}
does not define the FLRW anomaly abundance or the linear CMB kernel by itself.

\paragraph{Proof.}
The RAR law is written for a static weak-field acceleration vector sourced by a
resolved baryonic inhomogeneity.  An exactly homogeneous FLRW background has no
preferred spatial acceleration vector.  Its scalar data are background
densities and perturbation kernels, with no radial acceleration profile.  In
addition, the RAR map contains \(\sqrt{g_b}\), so its first derivative is
singular at \(g_b=0\) if treated as a homogeneous-background Taylor series.
Therefore the static law cannot be linearized into CMB equations without an
OPH parent functional.  Even a scalar parent functional must emit
\(\rho_A(a)\) and \(B_A(k,a)\), but that is still only the source-candidate
layer.  Physical Boltzmann status requires the finite covariant packet parent
above. \(\square\)

\subsection{Parent Collar Functional}

\paragraph{Scalar source target 0J.}
Let \(\mathcal C_x(a)\) be the finite collar family contributing to a coarse
FLRW cell \(x\) at scale factor \(a\).  Assume this finite-collar state
selection has a regular FLRW limit.  Let \(\omega[\rho_b]\) be the OPH collar
state selected by the baryonic source, and let
\begin{equation}
  R_C[\omega]=I_\omega(A:D|B).
\end{equation}
Define
\begin{equation}
  \rho_{A,\mathrm{eq}}(x,a)c^2
  =
  \frac{15}{8\pi^2\ell(a)^4}
  \int_{\mathcal C_x(a)}\dd\mu_C\,
  R_C[\omega[\rho_b]].
  \label{eq:parentcollar}
\end{equation}
Then \eqref{eq:parentcollar} is a scalar parent collar functional that emits a
background equilibrium source and a linear kernel:
\begin{equation}
  \bar\rho_{A,\mathrm{eq}}(a)=
  \rho_{A,\mathrm{eq}}[\bar\rho_b](a),
\end{equation}
\begin{equation}
  K_A^{(\rho)}(k,a)=
  \widehat{
  \left.
  \frac{\delta\rho_{A,\mathrm{eq}}(x,a)}
       {\delta\rho_b(x',a)}
  \right|_{\bar\rho_b}
  }(k),
  \qquad
  B_A(k,a)=\frac{\rho_b(a)}{\rho_A(a)}K_A^{(\rho)}(k,a).
\end{equation}

\paragraph{Conditional argument.}
The collar measure, the conditional mutual information, and the selected
collar state are quotient-local scalar data.  The integral is therefore a
scalar density on the coarse FLRW cell.  Evaluating the functional on a
homogeneous baryonic density gives the background equilibrium source.  Taking
the first functional derivative around that homogeneous state gives a
translation-invariant and rotation-invariant linear response kernel.  Its
Fourier transform is \(K_A^{(\rho)}(k,a)\), and division by the background density gives
\(B_A(k,a)\). \(\square\)

The finite-collar evaluation of \eqref{eq:parentcollar} is the mathematical
target for replacing a fitted dark-matter abundance, but it does not determine
the full fluid closure.  The scalar parent leaves open \(w_A\), \(c_{s,A}^2\),
\(\sigma_A\), \(Q_A^\mu\), \(\Gamma_{\mathrm{rec}}\), recipient stress, and
causal response.  Those data must come from the finite covariant packet parent
before the scalar tables can feed a physical Boltzmann module.

\paragraph{Proof obligation.}
Evaluate the collar measure and selected state from finite OPH screen/collar
data in the homogeneous, perturbative, galaxy, and cluster regimes; lift the
result to a finite covariant packet parent; prove exact stress-energy closure,
gauge independence of \(B_A\), causal/stable response, refinement convergence,
and CDM-limit recovery.

\subsection{Finite Repair Matrix}

\paragraph{Theorem 0K.}
Let \(\mathsf S\) be the finite packet-state space declared by one
fixed-cutoff overlap interface.  Let \(\pi_{\mathrm{eq}}(s|\rho_b)\) be the
equilibrium packet distribution emitted by the parent collar functional.  Let
\(q(s,s')\) be an inverse-closed local proposal rule generated by the declared
repair menu, with \(q(s,s')>0\) iff \(q(s',s)>0\).  Define, for \(s\ne s'\),
\begin{equation}
  K(s,s')
  =
  q(s,s')
  \min\!\left(
    1,
    \frac{\pi_{\mathrm{eq}}(s'|\rho_b)q(s',s)}
         {\pi_{\mathrm{eq}}(s|\rho_b)q(s,s')}
  \right),
\end{equation}
and
\begin{equation}
  K(s,s)=1-\sum_{s'\ne s}K(s,s').
\end{equation}
Then \(K\) is a finite stochastic repair transition matrix with stationary
distribution \(\pi_{\mathrm{eq}}\).  If the proposal graph is connected and
some rejection or hold probability is nonzero, \(K\) is irreducible and
aperiodic.

\paragraph{Proof.}
Each row sums to one by construction, so \(K\) is stochastic.  For
\(s\ne s'\), inverse-closed proposals make both directions available.  Hence
\begin{equation}
  \pi_{\mathrm{eq}}(s)K(s,s')
  =
  \min\!\left(
    \pi_{\mathrm{eq}}(s)q(s,s'),
    \pi_{\mathrm{eq}}(s')q(s',s)
  \right)
  =
  \pi_{\mathrm{eq}}(s')K(s',s).
\end{equation}
Detailed balance proves stationarity.  Connected proposal support gives
irreducibility, and a nonzero hold probability gives aperiodicity. \(\square\)

This construction declares \(K\) from OPH data: packet states, local proposals,
and the parent equilibrium distribution.

\paragraph{Scalar-load diagnostic.}
The executable packet diagnostic takes
\[
  \mathsf S=\{0,1,\ldots,40\},
  \qquad
  \pi_{\mathrm{eq}}(n)=
  \mathrm{Pois}\!\left(n;5.363470441\right)
\]
with truncation at \(40\), nearest-neighbor inverse-closed proposals, and
hold probability \(0.25\).  The resulting matrix has row-sum error
\(1.110\times10^{-16}\), detailed-balance error \(6.939\times10^{-18}\), and
\[
  \gamma_{\mathrm{rec}}=0.050717471.
\]
This diagnostic is a finite packet schema.  The OPH task is to derive that
schema, or its replacement, from the finite collar states and then attach the
stress, recipient, response, convergence, and frozen-likelihood receipts needed
for physical transfer.

\subsection{Finite Repair-Graph Gap}

\paragraph{Theorem 0L.}
For a declared fixed-cutoff repair graph with transition matrix \(K\), the
relaxation time used by clusters and perturbations is computable from the
second eigenvalue.  If \(K\) is reversible with stationary distribution
\(\pi_{\mathrm{eq}}\), then
\begin{equation}
  \gamma_{\mathrm{rec}}=1-\lambda_2(K)
\end{equation}
in discrete repair-step units, where \(\lambda_2\) is the largest nontrivial
eigenvalue in absolute value on the mean-zero subspace.  For continuous-time
generator \(L=K-\mathbf 1\), the gap is the smallest positive eigenvalue of
\(-L\).

\paragraph{Proof.}
Reversibility makes \(K\) self-adjoint in
\(L^2(\pi_{\mathrm{eq}})\).  The stationary vector has eigenvalue one.  Every
mean-zero mode decomposes into eigenvectors and decays by its eigenvalue per
repair step.  The slowest decaying nonstationary mode is therefore the largest
nontrivial eigenvalue, giving \(\gamma_{\mathrm{rec}}=1-\lambda_2(K)\).  The
continuous-time statement follows by exponentiating the generator. \(\square\)

This is an exact finite-matrix target for the repair/commit microphysics.  A
cluster timescale is a prediction only after this matrix or its OPH state
functional is supplied.

\subsection{Jacobi Repair Clock}

\paragraph{Theorem target 0M.}
Assume repair timing is set by the overlap-visible Jacobi mismatch rate of
transported information-defect records.  Let \(E_{ij}=C_{0i0j}\) be the
electric Weyl tensor in the local free-fall frame, equivalently the trace-free
Newtonian tidal tensor in the weak-field limit.  Homogeneous FLRW Ricci
focusing is excluded from this local cluster-offset clock.  The scalar repair
rate is
\begin{equation}
  \Gamma_J
  =
  \left(\frac{E_{ij}E^{ij}}{6}\right)^{1/4}.
\end{equation}
For a Newtonian point source this gives
\begin{equation}
  \Gamma_J=\sqrt{\frac{GM}{r^3}},
  \qquad
  \tau_J=\sqrt{\frac{r^3}{GM}}.
\end{equation}

\paragraph{Explicit premise.}
The repair-cycle clock is controlled by overlap-visible geodesic-deviation
mismatch of transported information-defect records.

\paragraph{Conditional argument.}
The relative acceleration of neighboring transported records is controlled by
the Jacobi operator.  The local cluster-offset part is the trace-free Weyl
electric tensor \(E_{ij}\); the homogeneous FLRW Ricci term changes the
background expansion and is excluded from the local offset between collisional
gas and transported record stress.  A scalar clock made from \(E_{ij}\), with no
chosen direction and no extra dimensional constant, must be a power of
\(E_{ij}E^{ij}\).  The fourth root has units of inverse time.  For a point
Newtonian source the eigenvalues of \(E_{ij}\) are
\(-2GM/r^3,GM/r^3,GM/r^3\) up to sign convention, hence
\(E_{ij}E^{ij}=6(GM/r^3)^2\).  The normalization in the theorem makes the
point-source result exact. \(\square\)

\paragraph{Proof obligation.}
Derive from OPH packet microphysics that the commit clock is the
overlap-visible Weyl-Jacobi mismatch clock, while pixel, de Sitter,
entropy-production, local-acceleration, and Ricci-curvature clocks are
excluded.  The
self-consistent prediction evaluates \(E_{ij}\) from the total weak-field
potential determined by baryons plus information-defect stress.

\paragraph{Failure consequence.}
Without that derivation, \(\Gamma_J\) is a no-fit dynamic diagnostic rather
than a repair-clock theorem.

\paragraph{Static-clock boundary.}
The static equilibrium equation contains \(\rho_{A,\mathrm{eq}}\) but no
relaxation rate.  Replacing \(\Gamma_{\mathrm{rec}}\) by any positive function
leaves the same settled solution.  The Jacobi clock is therefore a dynamic
OPH premise and has no theorem status inside the static RAR law.

\subsection{Homogeneous Anomaly State Selection}

\paragraph{Theorem target 0N.}
Under flat capacity-saturated FLRW state selection, the homogeneous anomaly
abundance is
\begin{equation}
  \Omega_{A0}
  =
  1-\Omega_{\Lambda,\OPH}
   -\Omega_{b0}-\Omega_{\nu0}-\Omega_{r0},
  \qquad
  \Omega_{\Lambda,\OPH}
  =
  \left(\frac{H_{\mathrm{dS}}}{H_0}\right)^2.
\end{equation}
If the homogeneous anomaly charge is transported as dust with no homogeneous
repair exchange, then
\begin{equation}
  \rho_A(a)=\rho_{A0}a^{-3}.
\end{equation}

\paragraph{Explicit premise.}
The homogeneous state is flat, capacity-saturated, and selected by a residual
Friedmann constraint after \(H_0\), baryons, neutrinos, and radiation are
supplied.

\paragraph{Conditional argument.}
The flat Friedmann equation at \(a=1\) is
\begin{equation}
  1
  =
  \Omega_{\Lambda,\OPH}
  +\Omega_{b0}+\Omega_{\nu0}+\Omega_{r0}+\Omega_{A0}.
\end{equation}
Solving for \(\Omega_{A0}\) gives the displayed residual.  Pressureless
homogeneous transport with zero homogeneous exchange obeys
\[
  \rho_A'+3\mathcal H\rho_A=0,
\]
whose solution is \(\rho_A(a)=\rho_{A0}a^{-3}\). \(\square\)

\paragraph{Proof obligation.}
Derive the conserved homogeneous information-defect charge \(Q_A\) from OPH
screen state data, or declare it as cosmological state data.

\paragraph{Failure consequence.}
The CMB dark abundance is conditional if \(Q_A\) is not selected by OPH
microphysics.  A nonzero \(\Gamma_{\mathrm{rec}}\) branch is also conditional
unless the same finite parent exposes the recipient stress that balances the
exchange.

\paragraph{Capacity boundary.}
The screen capacity fixes the dimensionless de Sitter capacity relation; the
SI de Sitter scale additionally uses the selected OPH scale certificate.  After
\(H_0\) and that scale display are supplied, it fixes \(\Omega_{\Lambda,\OPH}\).
It does not by itself fix the baryon density,
the neutrino density, the radiation density, or a conserved homogeneous
anomaly charge.  A standalone OPH cosmology needs a theorem of the form
\[
  Q_A=F_{\OPH}(\hbox{screen capacity},P,\hbox{screen state data})
\]
or an explicitly declared state-selection premise such as
\eqref{eq:flatresidual}.

\section{Falsifiers}

The OPH dark branch fails if any of the following happens:
\begin{enumerate}
\item OPH collar/cut microphysics cannot derive the activation law
\(p(x)=1-\exp[-\lambdac\sqrt{x}]\), or a replacement law that passes the same
galaxy tests.
\item The carried galaxy defect is not positive or not attractive.
\item The settled anomaly stress cannot give \(\Phi=\Psi\).
\item A full SPARC likelihood rejects both the unit branch and every
independently derived coefficient branch under the same nuisance treatment used
for comparison models.
\item Cluster lensing offsets require a relaxation law incompatible with old
settled-galaxy RAR.
\item CMB peaks, CMB lensing, BAO, weak lensing, or growth reject every frozen
finite covariant parent whose \(\Gamma_{\mathrm{rec}}\), \(B_A(k,a)\),
\(\rho_A(a)\), stress closure, recipient channel, and CDM-limit recovery remain
compatible with the galaxy branch.
\end{enumerate}

\section{Reproducibility}

The numerical tables were produced in the OPH cosmology workspace by:
\begin{verbatim}
python3 cosmology/scripts/dark_sector_measurement_ledger.py
python3 cosmology/scripts/dark_empirical_scorecard.py --quiet
python3 cosmology/scripts/sparc_systematic_likelihood.py
python3 cosmology/scripts/dark_perturbation_parameters.py
python3 cosmology/scripts/dark_repair_transition_matrix.py
python3 cosmology/scripts/dark_cluster_boltzmann_prelikelihood.py
python3 cosmology/scripts/dark_jacobi_repair_clock.py
python3 cosmology/scripts/dark_homogeneous_state_selection.py
python3 cosmology/scripts/dark_parent_collar_grid.py --diagnostic-scalar-load \
  --mu-eq 5.363470440729118 \
  --out cosmology/outputs/dark_parent_flat_capacity_state_selection.json
python3 cosmology/scripts/dark_cmb_bao_growth_s8_likelihood.py \
  --parent-grid cosmology/outputs/dark_parent_flat_capacity_state_selection.json
python3 cosmology/scripts/camb_fixed_neutrino_compare.py
\end{verbatim}
The SPARC data files used are public measurement tables from the SPARC page
\cite{sparcpage}.

\begin{thebibliography}{99}

\bibitem{ophcompact}
B. Mueller, A. Osika, M. Poneder, K. Xue, P. Nguyen, and M. A. Visser,
\emph{Recovering Relativity and the Standard Model from Observer Overlap Consistency},
OPH compact paper, release r1411, May 3, 2026.
\url{https://github.com/FloatingPragma/observer-patch-holography/blob/main/paper/recovering_relativity_and_standard_model_structure_from_observer_overlap_consistency_compact.pdf}

\bibitem{ophsynthesis}
B. Mueller, A. Osika, M. Poneder, K. Xue, B. Cassie, P. Nguyen, M. A. Visser, and K. A. Anirudha,
\emph{Observers Are All You Need}, OPH synthesis paper, release r1411,
May 3, 2026.
\url{https://github.com/FloatingPragma/observer-patch-holography/blob/main/paper/observers_are_all_you_need.pdf}

\bibitem{ophmicro}
B. Mueller, A. Osika, K. Xue, B. Cassie, and M. A. Visser,
\emph{Federated Echosahedral Screen Microphysics: Patch Hardware, Records, and Observer Synchronization in OPH},
OPH microphysics paper.
\url{https://github.com/FloatingPragma/observer-patch-holography/blob/main/paper/screen_microphysics_and_observer_synchronization.pdf}

\bibitem{ophparticles}
B. Mueller, A. Osika, M. Poneder, K. Xue, and M. A. Visser,
\emph{Deriving the Particle Zoo from Observer Consistency},
OPH particle paper, release r1411, May 3, 2026.
\url{https://github.com/FloatingPragma/observer-patch-holography/blob/main/paper/deriving_the_particle_zoo_from_observer_consistency.pdf}

\bibitem{planck2018}
Planck Collaboration,
\emph{Planck 2018 results. VI. Cosmological parameters},
arXiv:1807.06209.
\url{https://arxiv.org/abs/1807.06209}

\bibitem{lelli2016}
F. Lelli, S. S. McGaugh, and J. M. Schombert,
\emph{SPARC: Mass Models for 175 Disk Galaxies with Spitzer Photometry and Accurate Rotation Curves},
Astron. J. 152, 157, 2016, arXiv:1606.09251.
\url{https://arxiv.org/abs/1606.09251}

\bibitem{sparcpage}
F. Lelli, S. S. McGaugh, and J. M. Schombert,
SPARC public data page.
\url{https://astroweb.cwru.edu/SPARC/}

\bibitem{mcgaugh2016}
S. S. McGaugh, F. Lelli, and J. M. Schombert,
\emph{The Radial Acceleration Relation in Rotationally Supported Galaxies},
Phys. Rev. Lett. 117, 201101, 2016, arXiv:1609.05917.
\url{https://arxiv.org/abs/1609.05917}

\bibitem{li2018}
P. Li, F. Lelli, S. S. McGaugh, and J. M. Schombert,
\emph{Fitting the radial acceleration relation to individual SPARC galaxies},
Astron. Astrophys. 615, A3, 2018, arXiv:1803.00022.
\url{https://arxiv.org/abs/1803.00022}

\bibitem{milgrom1983}
M. Milgrom,
\emph{A modification of the Newtonian dynamics as a possible alternative to the hidden mass hypothesis},
Astrophys. J. 270, 365, 1983.
\url{https://ui.adsabs.harvard.edu/abs/1983ApJ...270..365M/abstract}

\bibitem{bekenstein2004}
J. D. Bekenstein,
\emph{Relativistic gravitation theory for the modified Newtonian dynamics paradigm},
Phys. Rev. D 70, 083509, 2004.
\url{https://doi.org/10.1103/PhysRevD.70.083509}

\bibitem{clowe2006}
D. Clowe, M. Bradac, A. H. Gonzalez, M. Markevitch, S. W. Randall,
C. Jones, and D. Zaritsky,
\emph{A direct empirical proof of the existence of dark matter},
Astrophys. J. Lett. 648, L109, 2006, arXiv:astro-ph/0608407.
\url{https://arxiv.org/abs/astro-ph/0608407}

\end{thebibliography}

\end{document}
